\documentclass[11pt]{amsart}

% Encoding and fonts
\usepackage[T1]{fontenc}
\usepackage[utf8]{inputenc}
\usepackage{lmodern}

% Math packages
\usepackage{amsmath}
\usepackage{amssymb}
\usepackage{amsfonts}
\usepackage{amsthm}
\usepackage{mathtools}

% Layout and utilities
\usepackage{geometry}
\usepackage{enumitem}
\usepackage[hidelinks]{hyperref}

\geometry{margin=1in}

\numberwithin{equation}{section}
\setcounter{tocdepth}{3}
\setcounter{secnumdepth}{3}

% ---------------------------------------------------------
% Theorem-like environments
% ---------------------------------------------------------

\theoremstyle{plain}
\newtheorem{theorem}{Theorem}[section]
\newtheorem{proposition}[theorem]{Proposition}
\newtheorem{lemma}[theorem]{Lemma}
\newtheorem{corollary}[theorem]{Corollary}
\newtheorem{conjecture}[theorem]{Conjecture}

\theoremstyle{definition}
\newtheorem{definition}[theorem]{Definition}
\newtheorem{experiment}[theorem]{Experiment}
\newtheorem{heuristic}[theorem]{Heuristic}
\newtheorem{question}[theorem]{Question}
\newtheorem{hypothesis}[theorem]{Hypothesis}

\theoremstyle{remark}
\newtheorem{remark}[theorem]{Remark}
\newtheorem{warning}[theorem]{Warning}

% ---------------------------------------------------------
% Notation
% ---------------------------------------------------------

\DeclareMathOperator{\sgn}{sgn}
\DeclareMathOperator{\li}{li}
\newcommand{\Ms}{M_{\mathrm{s}}}
\newcommand{\Msk}{M_{\mathrm{sk}}}
\newcommand{\Ap}{\mathcal{A}_{+}}
\newcommand{\Am}{\mathcal{A}_{-}}
\newcommand{\Anet}{\mathcal{A}}
\newcommand{\eps}{\varepsilon}
\newcommand{\Pp}{P^{+}}
\newcommand{\Cc}{\mathcal{C}}
\newcommand{\Sq}{\mathcal{S}}
\newcommand{\En}{\mathcal{E}}
\newcommand{\mt}{\widetilde{m}}

\title[The area function of the Mertens excursion skeleton]{The area function of the Mertens excursion skeleton\\
(research notebook, version 9 --- revised merged edition)}

\author{V.~Casa\~na-Giner}
\date{25 September 2026}

\begin{document}

\maketitle

\tableofcontents

\bigskip

\begin{remark}
\label{rem:exploratory_nature}
This notebook is exploratory. Its purpose is to study the Mertens function
through its excursions between zero runs, and to separate what is
forced by the local geometry of those excursions from what requires
arithmetic input. It contains exact identities, conditional statements,
heuristics and archived statements from earlier versions. Every statement
is to be read with an explicit evidentiary status. Version 9 separates
\emph{[PUBLISHED INPUT]}, \emph{[PROVED HERE]}, \emph{[CONDITIONAL]}, and
\emph{[HEURISTIC]}.  A rigorous estimate that is too weak for the stated RH
objective is additionally marked \emph{[PROVED BUT WEAK]}; when the present
argument yields no estimate of the required type the text says
\emph{[NO USEFUL BOUND]}. Earlier statements that are false or superseded are
retained only when their failure is informative, and are marked
\emph{[FALSE / WITHDRAWN]}.
\end{remark}

\section{Introduction and conventions}
\label{sec:introduction}

The main object is the Mertens function $M$, decomposed into excursions
between zero runs, together with three area curves built from a
regularized skeleton of those excursions
(Section~\ref{sec:definitions}). Sections~\ref{sec:formulas}
to~\ref{sec:criteria} record exact formulas, sharp bounds and criteria
equivalent to the Riemann hypothesis. Section~\ref{sec:lacunary} contains
the two new geometric lacunarity bounds for long excursions, the stronger
published Kaczorowski--Pintz multiplicative localization of actual Mertens
returns, and the resulting localized RH criteria, including the
square-truncated criterion: in each excursion it is enough to
control the range $x_0<t\le\min\{x_1,x_0^2\}$.  Equivalently, one may ignore
the part of every square-exceptional excursion lying beyond $x_0^2$.
Sections~\ref{sec:prime_horizon} to~\ref{sec:local_large_prime} then examine
the square trial-division horizon and large-prime decompositions.

\medskip
\noindent\textbf{Status convention.} The following labels are literal.
\emph{[PUBLISHED INPUT]} means a theorem taken from the cited literature.
\emph{[PROVED HERE]} means an unconditional statement whose proof is supplied
in this notebook. \emph{[CONDITIONAL]} names its hypothesis explicitly.
\emph{[HEURISTIC]} is unproved. \emph{[PROVED BUT WEAK]} means rigorous but
insufficient for the stated RH objective. \emph{[NO USEFUL BOUND]} means
that the present argument supplies no non-trivial estimate of the required
type. \emph{[FALSE / WITHDRAWN]} identifies an earlier inference that is not
valid.  An unconditional equivalence with RH is a theorem and does not assume
RH.

\medskip
\noindent\textbf{Status lines.} Each section opens with a status line that
separates standard material, literature-based material, original
observations, heuristics and conjectures, as required by the authoring rules
of this notebook.

\medskip
\noindent\textbf{Notation changes with respect to version 5.} The symbol
$m(x)$ of version 5 is replaced by $M(x)$. The unregularized skeleton, called
$L$ in version 5, is denoted $\Msk$. The squarefree counting function, called
$Q(y)$ in version 5, is denoted $\Sq(y)$, since $Q_\nu$ denotes the quadratic
energy of an excursion. The Dirichlet energy, called $E(x)$ in version 5, is
denoted $\En(x)$. The coefficient called $A(R)$ in version 5 is denoted
$\Theta(R)$, since $\Anet$ denotes the black curve. The ratio $t/u^2$ keeps
the name $R$; the residual $M-\Msk$ of version 5, also called $R$ there, is
written out explicitly. The defect parameter of
Remark~\ref{rem:conditioning} is denoted $d_\nu$.


\section{Excursions and the regularized skeleton}
\label{sec:definitions}

\noindent\emph{Status.} The segmentation of $M$ by zero runs and the record
of its extrema are standard computational objects; Hurst~\cite{Hurst18}
computed $M(x)$ for all $x\le10^{16}$, recording all zeros and extrema.  The
area and quadratic functionals of lattice excursions are classical
\cite{Takacs91,NguyenThe04,Richard05,Janson07,DKW15,CP17}.  What is specific
to this notebook is the regularized triangular skeleton and its use on the
deterministic M\"obius walk.  All statements in this section are
unconditional.

\medskip

Let $\mu$ be the M\"obius function and
\begin{equation}
\label{eq:mertens_def}
M(x)=\sum_{n\le x}\mu(n),\qquad x\ge0,
\end{equation}
the Mertens function, piecewise constant and right-continuous, with jumps
$\mu(n)\in\{-1,0,1\}$ at the integers.

\begin{definition}[Excursions and zero runs]
\label{def:excursion}
An \emph{excursion} is an interval
\[
I_\nu=(x_{0,\nu},x_{1,\nu}]=(u_\nu,v_\nu]
\]
with $M(x_{0,\nu})=M(x_{1,\nu})=0$,
$M(x_{0,\nu}+1)=\sigma_\nu\in\{\pm1\}$, and
$\sigma_\nu M(n)\ge1$ for every integer
$x_{0,\nu}<n<x_{1,\nu}$. Thus $x_{1,\nu}$ is the first return to zero after
the departure from $x_{0,\nu}$. Consecutive excursions may be separated by
a \emph{zero run}
\[
[x_{1,\nu},x_{0,\nu+1}],
\]
a maximal set of consecutive integers on which $M=0$; it may consist of a
single point. Zero runs are not excursions. The first excursion is $(0,2]$.
We retain $u_\nu,v_\nu$ as legacy aliases for $x_{0,\nu},x_{1,\nu}$ in
formulas inherited from versions 5--6.

Put
\[
\ell_\nu=x_{1,\nu}-x_{0,\nu},\qquad
H_\nu=\max_{x_{0,\nu}<n<x_{1,\nu}}|M(n)|,
\]
\[
\tau_\nu=\min\{n\in I_\nu:|M(n)|=H_\nu\},\qquad
a_\nu=\tau_\nu-x_{0,\nu},\qquad b_\nu=x_{1,\nu}-\tau_\nu.
\]
Every excursion has $H_\nu\ge1$ and $\ell_\nu\ge2$.
\end{definition}

\begin{lemma}[Height and step count]
\label{lem:height_steps}
For every excursion,
\begin{equation}
\label{eq:height_steps}
H_\nu\le\min(a_\nu,b_\nu)\le\frac{\ell_\nu}{2}.
\end{equation}
\end{lemma}

\begin{proof}
The steps of $M$ are $0$ or $\pm1$, so at least $H_\nu$ steps are needed to
reach height $H_\nu$ from zero and at least $H_\nu$ more to return to zero.
\end{proof}

\begin{definition}[Crossing and touching zero runs]
\label{def:crossing_touching}
The zero run $[x_{1,\nu},x_{0,\nu+1}]$ between consecutive excursions is a
\emph{crossing run} if $\sigma_{\nu+1}=-\sigma_\nu$ and a \emph{touching run}
if $\sigma_{\nu+1}=\sigma_\nu$.
\end{definition}

\begin{definition}[Regularized skeleton]
\label{def:Ms}
The \emph{skeleton} $\Ms$ is the continuous, piecewise linear function
defined on every excursion by
\begin{equation}
\label{eq:Ms_def}
\Ms(t)=
\begin{cases}
\sigma_\nu\sqrt{H_\nu}\dfrac{t-x_{0,\nu}}{a_\nu},
 &x_{0,\nu}\le t\le\tau_\nu,\\[2ex]
\sigma_\nu\sqrt{H_\nu}\dfrac{x_{1,\nu}-t}{b_\nu},
 &\tau_\nu\le t\le x_{1,\nu},
\end{cases}
\end{equation}
and by $\Ms\equiv0$ on zero runs.  The \emph{unregularized skeleton} $\Msk$
is defined in the same way with peak $\sigma_\nu H_\nu$.
\end{definition}

\begin{definition}[Terminal block of a truncated computation]
\label{def:Ms-terminal}
Let $N\ge2$ and let $u=\max\{n\le N:M(n)=0\}$.  If $u=N$ there is no
terminal block. Otherwise put
\[
\sigma_{\rm term}=\sgn M(N),\qquad
H_{\rm term}=\max_{u<n\le N}|M(n)|,
\]
\[
\tau_{\rm term}=\min\{n\in(u,N]:|M(n)|=H_{\rm term}\}.
\]
On $[u,N]$ the regularized skeleton is the segment from $(u,0)$ to
$(\tau_{\rm term},\sigma_{\rm term}\sqrt{H_{\rm term}})$ followed, if
$\tau_{\rm term}<N$, by the segment ending at
$(N,\sigma_{\rm term}\sqrt{|M(N)|})$. Its signed area is
\[
\Delta\Anet(u,N]=\frac{\sigma_{\rm term}}2
\Bigl[\sqrt{H_{\rm term}}(N-u)+
\sqrt{|M(N)|}(N-\tau_{\rm term})\Bigr].
\]
The terminal skeleton is provisional: after the next return to zero the
completed excursion acquires a different descending side.  Independently of
this provisional skeleton, for a real truncation point $x$ with last zero
$u(x)\le\lfloor x\rfloor$ define the actual terminal quantities
\[
\beta_{\rm term}(x):=\sum_{u(x)<n\le\lfloor x\rfloor}|M(n)|,\qquad
Q_{\rm term}(x):=\sum_{u(x)<n\le\lfloor x\rfloor}M(n)^2,
\]
and $L_x:=\lfloor x\rfloor-u(x)$.  These depend only on $M$, not on the
provisional skeleton.
\end{definition}

\begin{definition}[The three curves]
\label{def:areas}
With $\Ms^{\pm}=\max(\pm\Ms,0)$,
\begin{equation}
\label{eq:areas_def}
\Ap(x)=\int_0^x\Ms^+(t)\,dt,\qquad
\Am(x)=-\int_0^x\Ms^-(t)\,dt,\qquad
\Anet(x)=\int_0^x\Ms(t)\,dt.
\end{equation}
$\Ap$ is nondecreasing, $\Am$ is nonincreasing, and both are constant on
zero runs and on excursions of the opposite sign.
\end{definition}

\begin{remark}[The family of $g$-skeletons]
\label{rem:g-skeleton}
For any $g:\mathbb N_0\to[0,\infty)$ with $g(0)=0$ and $g(h)>0$ for $h\ge1$,
one may replace the apex height $\sqrt H$ by $g(H)$. The signed area increment
becomes $\tfrac12\sigma_\nu g(H_\nu)\ell_\nu$.  The excursion correlation identities introduced below concern the M\"obius
increments and are independent of the choice of $g$.  The choice $g(H)=\sqrt H$ is a normalization,
not a geometric necessity.
\end{remark}

\section{Exact formulas}
\label{sec:formulas}

\noindent\emph{Status.} Exact elementary formulas with complete proofs.  The
critical-geometry lemma includes the correction for positive-length zero
plateaus.

\begin{lemma}[Increment per excursion]
\label{lem:jump}
For every complete excursion,
\begin{equation}
\label{eq:jump}
\Anet(x_{1,\nu})-\Anet(x_{0,\nu})
=\sigma_\nu\frac{\sqrt{H_\nu}\,\ell_\nu}{2}.
\end{equation}
The discrete sum agrees:
\begin{equation}
\label{eq:jump_discrete}
\sum_{n\in I_\nu}\Ms(n)
=\sigma_\nu\sqrt{H_\nu}\frac{\ell_\nu}{2}.
\end{equation}
\end{lemma}

\begin{proof}
The two pieces are triangles of height $\sqrt{H_\nu}$ and bases $a_\nu,b_\nu$.
For the ascent, the integer values are indexed by $j=1,\ldots,a_\nu$ and
sum to
\[
\sqrt{H_\nu}\sum_{j=1}^{a_\nu}\frac{j}{a_\nu}
=\sqrt{H_\nu}\frac{a_\nu+1}{2}.
\]
For the descent the apex is excluded to avoid double counting; with
$r=1,\ldots,b_\nu$ the values sum to
\[
\sqrt{H_\nu}\sum_{r=1}^{b_\nu}\frac{b_\nu-r}{b_\nu}
=\sqrt{H_\nu}\frac{b_\nu-1}{2}.
\]
Their sum is $\sqrt{H_\nu}\ell_\nu/2$.
\end{proof}

\begin{proposition}[Explicit form of the three curves]
\label{prop:explicit}
Let $x\in I_\nu$. Then
\begin{equation}
\label{eq:explicit_A}
\Anet(x)=
\begin{cases}
\Anet(x_{0,\nu})+\sigma_\nu\sqrt{H_\nu}
\dfrac{(x-x_{0,\nu})^2}{2a_\nu},&x\le\tau_\nu,\\[2ex]
\Anet(x_{0,\nu})+\sigma_\nu\sqrt{H_\nu}
\left(\dfrac{a_\nu}{2}+(x-\tau_\nu)-
\dfrac{(x-\tau_\nu)^2}{2b_\nu}\right),&x\ge\tau_\nu.
\end{cases}
\end{equation}
At excursion endpoints,
\begin{align}
\Anet(x_{1,k})&=\sum_{\nu\le k}\sigma_\nu\frac{\sqrt{H_\nu}\ell_\nu}{2},\label{eq:explicit_net}\\
\Ap(x_{1,k})&=\sum_{\substack{\nu\le k\\\sigma_\nu=+1}}
\frac{\sqrt{H_\nu}\ell_\nu}{2},\label{eq:explicit_plus}\\
\Am(x_{1,k})&=-\sum_{\substack{\nu\le k\\\sigma_\nu=-1}}
\frac{\sqrt{H_\nu}\ell_\nu}{2}.\label{eq:explicit_minus}
\end{align}
The curves are piecewise quadratic and $C^1$.
\end{proposition}

\begin{proof}
Integrate \eqref{eq:Ms_def}; the endpoint formulas follow from
Lemma~\ref{lem:jump}.
\end{proof}

\begin{definition}[Strict inflection used in this notebook]
A point $c$ is called a \emph{strict inflection point} if there exists
$\delta>0$ such that the function is strictly convex on one side of $c$ and
strictly concave on the other.  This definition does not require the second
derivative to exist at $c$.
\end{definition}

\begin{lemma}[Critical geometry of $\Anet$]
\label{lem:critical}
Let $R=[p,q]$ be the zero run between an excursion of sign $\sigma_-$ ending
at $p$ and the next excursion of sign $\sigma_+$ starting at $q$.
\begin{enumerate}[label=\textup{(\roman*)}]
\item Every apex $\tau_\nu$ is a strict inflection point.
\item If $\sigma_+=-\sigma_-$, then $R$ is an extremal plateau; it is a
strict local extremum exactly when $p=q$.
\item If $\sigma_+=\sigma_-=:\sigma$, then $\Anet$ has no local extremum on
$R$.  Its concavity has sign $-\sigma$ immediately to the left and
$+\sigma$ immediately to the right.  If $p=q$, the point is a strict
inflection; if $p<q$, no point of $R$ is a strict inflection.
\item Strict local extrema occur exactly at isolated crossing zeros, and
plateaus of extrema exactly at crossing runs of positive length.
\end{enumerate}
\end{lemma}

\begin{proof}
On the descending side of the left excursion,
$\sgn\Anet'=\sigma_-$ and $\sgn\Anet''=-\sigma_-$.  On the ascending side
of the right excursion, $\sgn\Anet'=\sigma_+$ and
$\sgn\Anet''=\sigma_+$.  On $R$, $\Anet'=0$.  The assertions follow, and at
every apex the sign of $\Anet''$ reverses.
\end{proof}

\begin{remark}[Status of the version 5 statement]
\label{rem:critical_v5}
Version~5 stated that every touching zero is an inflection point.  This is
correct verbatim only for a touching run consisting of one zero.  On a
positive-length touching run, $\Anet''$ has the pattern
$-\sigma,0,+\sigma$ and the change of concavity is carried by the whole
plateau.
\end{remark}

\section{Internal correlations of an excursion}
\label{sec:internal_corr}

\noindent\emph{Status.} Exact elementary identities with complete proofs.
The triangular identity holds for arbitrary real sequences; the fixed-length
upper bound for $Q$ is sharp with leading constant $1/12$, correcting the
constant $1/4$ in version~6.

For a real block $x_1,\dots,x_L$, put
\begin{equation}
\label{eq:block_def}
S_t=\sum_{j=1}^t x_j,\qquad
C_h=\sum_{j=1}^{L-h}x_jx_{j+h}.
\end{equation}

\begin{lemma}[Triangular identity]
\label{lem:triangular-corr}
For every real finite sequence,
\begin{equation}
\label{eq:general-triangular}
\sum_{t=1}^{L-1}S_t^2+\sum_{h=1}^{L-1}hC_h
=S_L\sum_{t=1}^{L-1}S_t.
\end{equation}
\end{lemma}

\begin{proof}
Expanding the partial sums gives
\[
\sum_{t=1}^{L}S_t^2
=\sum_{i,j=1}^{L}x_ix_j\bigl(L+1-\max(i,j)\bigr).
\]
Use
$\max(i,j)=\tfrac12(i+j+|i-j|)$.  The part containing $|i-j|$ is
\[
\frac12\sum_{i,j}x_ix_j|i-j|
=\sum_{h=1}^{L-1}h\sum_{j=1}^{L-h}x_jx_{j+h}
=\sum_{h=1}^{L-1}hC_h.
\]
For the remaining terms use
$\sum_i i x_i=LS_L-\sum_{t=1}^{L-1}S_t$.  After simplifying and subtracting
$S_L^2$ from both sides one obtains \eqref{eq:general-triangular}.
\end{proof}

\begin{corollary}[Bridge case]
\label{cor:bridge}
If $S_L=0$, then
\begin{equation}
\label{eq:bridge-triangular}
\sum_{t=1}^{L-1}S_t^2=-\sum_{h=1}^{L-1}hC_h.
\end{equation}
\end{corollary}

\begin{remark}[Hermitian form]
\label{rem:hermitian}
For complex $x_j$, with
$C_h=\sum_{j=1}^{L-h}x_j\overline{x_{j+h}}$,
\[
\sum_{t=1}^{L-1}|S_t|^2+\sum_{h=1}^{L-1}h\Re C_h
=\Re\!\left(\overline{S_L}\sum_{t=1}^{L-1}S_t\right).
\]
\end{remark}

For a complete excursion, orient the increments by
\begin{equation}
\label{eq:orientation}
x_j=\sigma_\nu\mu(x_{0,\nu}+j),\qquad
S_t=|M(x_{0,\nu}+t)|.
\end{equation}
Then $S_t>0$ for $1\le t<\ell_\nu$ and $S_{\ell_\nu}=0$. Put
\begin{equation}
\label{eq:Q_def}
Q_\nu=\sum_{t=1}^{\ell_\nu-1}S_t^2,\qquad
C_h^{(\nu)}=\sum_{j=1}^{\ell_\nu-h}x_jx_{j+h}.
\end{equation}
Thus
\begin{equation}
\label{eq:Qcorr}
Q_\nu=-\sum_{h=1}^{\ell_\nu-1}hC_h^{(\nu)}.
\end{equation}

\begin{lemma}[Zeroth moment]
\label{lem:zeroth-corr}
If $q_\nu=\sum_jx_j^2$, then
\begin{equation}
\label{eq:zeroth-corr}
\sum_{h=1}^{\ell_\nu-1}C_h^{(\nu)}=-\frac{q_\nu}{2}.
\end{equation}
\end{lemma}

\begin{proof}
Since $S_{\ell_\nu}=\sum_jx_j=0$,
\[
0=\left(\sum_jx_j\right)^2
=\sum_jx_j^2+2\sum_{h=1}^{\ell_\nu-1}C_h^{(\nu)}
=q_\nu+2\sum_{h=1}^{\ell_\nu-1}C_h^{(\nu)}.
\]
Rearranging gives the claim.
\end{proof}

\begin{lemma}[Extreme lag]
\label{lem:extreme-lag}
For every complete excursion,
\[
C_{\ell_\nu-1}^{(\nu)}=-1,
\]
and $q_\nu$ is even and at least $2$.
\end{lemma}

\begin{proof}
The first oriented step is $+1$.  Immediately before the first return to
zero the path is at height $1$, so the last step is $-1$.  The bridge
condition forces equally many positive and negative active steps.
\end{proof}

\begin{proposition}[Quadratic decomposition by excursions]
\label{prop:quadratic-decomp}
For real $x\ge1$, let
\[
u(x)=\max\{n\le\lfloor x\rfloor:M(n)=0\},\qquad
L_x=\lfloor x\rfloor-u(x).
\]
For $1\le j\le L_x$ put $y_j=\mu(u(x)+j)$,
$S_t=M(u(x)+t)$ and
$C_h^{\rm term}=\sum_{j=1}^{L_x-h}y_jy_{j+h}$. Then
\begin{equation}
\label{eq:quadratic-decomp-general}
\sum_{n\le x}M(n)^2
=\sum_{\nu:x_{1,\nu}\le u(x)}Q_\nu+Q_{\rm term}(x),
\end{equation}
where
\begin{equation}
\label{eq:terminal-Q-exact}
Q_{\rm term}(x)=M(\lfloor x\rfloor)^2
+M(\lfloor x\rfloor)\sum_{t=1}^{L_x-1}M(u(x)+t)
-\sum_{h=1}^{L_x-1}hC_h^{\rm term}.
\end{equation}
If $L_x=0$, $Q_{\rm term}(x)=0$.  At a zero $x$,
\begin{equation}
\label{eq:quadratic-decomp-zero}
\sum_{n\le x}M(n)^2
=\sum_{x_{1,\nu}\le x}Q_\nu
=-\sum_{x_{1,\nu}\le x}\sum_hhC_h^{(\nu)}.
\end{equation}
\end{proposition}

\begin{proof}
Every integer up to $u(x)$ lies either in a zero run or in a unique excursion
interior.  Apply Lemma~\ref{lem:triangular-corr} to the terminal block and
add its final $S_{L_x}^2$.
\end{proof}

\begin{remark}[Blindness to pairs crossing excursion boundaries]
\label{rem:blindness-quadratic}
Globally,
\[
\sum_{n\le x}M(n)^2
=\sum_{a,b\le x}\mu(a)\mu(b)(x+1-\max(a,b)).
\]
To see the cancellation precisely, first fix one index outside a complete
excursion and sum the other index over that excursion.  The weight
$x+1-\max(a,b)$ is affine in the index that ranges across the block on each
side of the fixed index, while the endpoint bridge relation
$\sum_{n\in I_\nu}\mu(n)=0$ removes the constant cross-block contribution.
Iterating this separation over completed excursions leaves the within-block
terms encoded by Proposition~\ref{prop:quadratic-decomp} and the unique
unfinished terminal block.  The proposition, rather than this sentence,
is the exact identity being used.
\end{remark}

Let
\begin{equation}
\label{eq:beta_def}
\beta_\nu=\sum_{t=1}^{\ell_\nu-1}S_t
=\sum_{n\in I_\nu}|M(n)|.
\end{equation}
Cauchy--Schwarz gives the sharp factor
\begin{equation}
\label{eq:beta-Q}
\beta_\nu^2\le(\ell_\nu-1)Q_\nu,
\end{equation}
with equality exactly for a flat height-$1$ excursion.

To separate support from sign geometry, let
$P_h^{(\nu)}=\#\{j:x_jx_{j+h}\ne0\}$ and let $D_h^{(\nu)}$ count the
opposite-sign pairs among them. Then
\begin{equation}
\label{eq:C-P-D}
C_h^{(\nu)}=P_h^{(\nu)}-2D_h^{(\nu)},
\end{equation}
and
\begin{equation}
\label{eq:Q-P-D}
Q_\nu=\sum_{h=1}^{\ell_\nu-1}h(2D_h^{(\nu)}-P_h^{(\nu)}).
\end{equation}

\begin{proposition}[Sharp bounds at fixed length and height]
\label{prop:sharp-LH}\label{rem:height-Q}
Let $S_0=S_L=0$, $S_t\ge1$ for $0<t<L$,
$S_t-S_{t-1}\in\{-1,0,1\}$ and $H=\max S_t$. Then $L\ge2H$ and
\begin{align}
H^2+L-2H&\le\beta\le H(L-H),\label{eq:beta-LH}\\
\frac{2H^3+H}{3}+L-2H&\le Q\le
\frac{2H^3+H}{3}+(L-2H)H^2.\label{eq:Q-LH}
\end{align}
All four bounds are attained. In particular
\begin{equation}
\label{eq:height_Q_min}
Q\ge\frac{2H^3+H}{3},
\end{equation}
and $Q\ge L-1$.
\end{proposition}

\begin{proof}
Before the first visit to $H$ and after it the path must visit each level
$1,\dots,H-1$.  Together with one visit to $H$, these $2H-1$ designated interior times
contribute
\[
2(1+\cdots+(H-1))+H=H^2
\]
to $\beta$, and
\[
2\sum_{j=1}^{H-1}j^2+H^2=\frac{2H^3+H}{3}
\]
to $Q$.
The remaining $L-2H$ times lie between heights $1$ and $H$. Keeping them all
at height $1$ attains the lower bounds; keeping them at height $H$ attains
the upper bounds.
\end{proof}

\begin{corollary}[Sharp fixed-length bounds]
\label{cor:fixed-L}
For $L\ge2$,
\[
\beta\le\left\lfloor\frac{L^2}{4}\right\rfloor,
\]
and
\[
Q\le
\begin{cases}
\dfrac{L(L^2+2)}{12},&L\text{ even},\\[1ex]
\dfrac{L(L^2-1)}{12},&L\text{ odd}.
\end{cases}
\]
Both are attained by the tent path, with one horizontal step at maximal height when $L$
is odd. In particular
\begin{equation}
\label{eq:Q_upper}
Q\le\frac{L^3+2L}{12}.
\end{equation}
\end{corollary}

\begin{proof}
Use $S_t\le\min(t,L-t)$ and sum this bound and its square.
\end{proof}

\section{Where \texorpdfstring{$\zeta$}{zeta} enters}
\label{sec:zeta}

\noindent\textbf{Status: MIXED.} The residue formula for $M$ is standard
\cite[\S14.27]{Titchmarsh1986}, with the usual summation convention and
multiplicity caveats. The displayed non-truncated expansion for
$\widetilde A$ below is \emph{FORMAL / NOT PROVED UNDER THE STATED
HYPOTHESES}; it is retained only as a residue calculation. A rigorous use
requires a truncated contour formula with an explicit remainder, or an
additional hypothesis controlling negative moments of $\zeta'(\rho)$.
Proposition~\ref{prop:dipoles} is an exact elementary identity proved here.

\medskip

The black curve depends on $M$ nonlinearly, through $\sqrt{H_\nu}$ and
through the selection of $u_\nu,\tau_\nu,v_\nu$. No closed expression of its
Mellin transform in terms of $\zeta$ is known to the author (version 5
asserted that none exists; that assertion is unproved and is withdrawn). The nodes $u_\nu,v_\nu$ are zeros of $M$. If the nontrivial zeros are simple, a standard residue formula (with no RH
assumption) is obtained only after choosing an admissible sequence of
truncation heights $T_k\to\infty$; the symbol $\sum_\rho$ below always means
that limiting procedure, not an assertion of ordinary or absolute
convergence:
\begin{equation}
\label{eq:explicitM}
  M(x)=\sum_{\rho}\frac{x^\rho}{\rho\,\zeta'(\rho)}-2
  +\sum_{n\ge1}\frac{(-1)^{n-1}(2\pi/x)^{2n}}{(2n)!\,n\,\zeta(2n+1)},
\end{equation}
for non-integer $x$, the sum over $\rho$ taken along a suitable sequence
$|\gamma|<T_k\to\infty$ \cite[\S14.27]{Titchmarsh1986}. For
$\widetilde{A}(x):=\int_0^xM(t)\,dt$ one has, for $\operatorname{Re}s>1$,
\begin{equation}
\label{eq:mellin_A}
  \int_1^\infty \widetilde{A}(x)x^{-s-2}\,dx=\frac{1}{s(s+1)\zeta(s)}.
\end{equation}
The corresponding inverse-Mellin residue computation is
\begin{equation}
\label{eq:explicitA}
  \widetilde{A}(x)\ \sim_{\rm formal}\
  \sum_\rho\frac{x^{\rho+1}}{\rho(\rho+1)\zeta'(\rho)}-2x+12
  +\text{(residues from the trivial zeros)}.
\end{equation}
Here the two displayed elementary residues are not placeholders: at $s=0$
we have $x/\zeta(0)=-2x$, while at $s=-1$ the residue is
$1/[(-1)\zeta(-1)]=12$.
\textbf{Status: FORMAL, NOT A THEOREM AS WRITTEN.} No untruncated equality
or $O(x^{-1})$ remainder is asserted here. The factor $(\rho+1)^{-1}$
improves convergence, but RH plus simplicity does not by itself control
$1/\zeta'(\rho)$ strongly enough. Rigorous untruncated statements require
additional input of Gonek--Hejhal type; see \cite{Ng2004}. The area of $\Msk$ differs from
$\widetilde{A}$ by explicit local terms:

\begin{proposition}[Dipole identity]
\label{prop:dipoles}
\emph{[PROVED HERE]} For each excursion let its \emph{legs} be the integer
sets $\lambda=(u_\nu,\tau_\nu]\cap\mathbb{Z}$ and $(\tau_\nu,v_\nu]\cap\mathbb{Z}$,
with midpoints $c_\lambda=(u_\nu+\tau_\nu)/2$ and $(\tau_\nu+v_\nu)/2$. Then
for every zero $v_k$,
\begin{equation}
\label{eq:dipoles}
  \int_0^{v_k}\Msk(t)\,dt=\widetilde{A}(v_k)
  +\sum_{\lambda\subset(0,v_k]}\ \sum_{m\in\lambda}\mu(m)\,(m-c_\lambda).
\end{equation}
\end{proposition}

\begin{proof}
On an ascent leg, $\int_{u}^{\tau}\Msk=\sigma H a/2$ and
$\int_u^\tau M=\sum_{m=u+1}^{\tau}\mu(m)(\tau-m)$, while
$\sigma H=\sum_{m=u+1}^{\tau}\mu(m)$; the difference is
$\sum_m\mu(m)(a/2-\tau+m)=\sum_m\mu(m)(m-c_\lambda)$. On a descent leg,
$\int_\tau^v\Msk=\sigma Hb/2$ and
$\int_\tau^vM=b\,\sigma H+\sum_{m=\tau+1}^{v}\mu(m)(v-m)$ with
$\sum_{m=\tau+1}^{v}\mu(m)=-\sigma H$; the difference is
$\sum_m\mu(m)(m-(v-b/2))$; since
$v-b/2=(\tau+v)/2=c_\lambda$, this is the required centered dipole.  Sum
over legs.
\end{proof}

\section{Bounds}
\label{sec:bounds}

\noindent\textbf{Status: [PUBLISHED INPUT] + [PROVED HERE] + one [FALSE / WITHDRAWN] inference.} The
transfer argument is elementary; the input bounds for $M$ are standard
\cite{Walfisz1963,OdlyzkoTeRiele1985,KotnikTeRiele2006,Hurst18}. The
former inference from an excursion-height bound indexed by $v_\nu$ to a
pointwise zero-free half-plane is \emph{FALSE / WITHDRAWN}; see
Remark~\ref{rem:onedirection}.

\begin{lemma}[Transfer]
\label{lem:transfer}
\emph{[PROVED HERE]} Let $x\ge1$.  In the infinite-function case assume a
zero $x^*\ge x$ has been selected (for example the first such zero, if it
exists).  In a finite computation truncated at $N\ge x$, replace $x^*$ by
$N$ and use the provisional terminal skeleton of
Definition~\ref{def:Ms-terminal}.  If
$|M(n)|\le c\,n^{1/2}$ for all $1\le n\le x^*$ (so necessarily $c\ge1$),
then
\begin{equation}
\label{eq:transfer_bound}
  |\Ms(t)|\le c^{1/2}t^{1/4}\quad(0\le t\le x),\qquad
  \Ap(x)+|\Am(x)|\le\tfrac45\,c^{1/2}x^{5/4}.
\end{equation}
\end{lemma}

\begin{proof}
For $t\ge\tau_\nu$, $|\Ms(t)|\le\sqrt{H_\nu}\le c^{1/2}\tau_\nu^{1/4}\le
c^{1/2}t^{1/4}$. For $u_\nu\le t\le\tau_\nu$ the chord from $(u_\nu,0)$ to
$(\tau_\nu,\sqrt{H_\nu})$ lies below the concave function $c^{1/2}t^{1/4}$,
which is $\ge0$ at $u_\nu$ and $\ge\sqrt{H_\nu}$ at $\tau_\nu\le x^*$.
Integrate.
\end{proof}

\begin{corollary}
\label{cor:uncond}
\emph{[PUBLISHED INPUT + PROVED HERE]} From the Vinogradov--Korobov bound
\cite{Walfisz1963}
\begin{equation}
\label{eq:VK_M}
  M(x)\ll x\exp\bigl(-c(\log x)^{3/5}(\log\log x)^{-1/5}\bigr)
\end{equation}
one obtains, for some $c'>0$,
\begin{equation}
\label{eq:uncond_area}
  \Ap(x)+|\Am(x)|\ll x^{3/2}\exp\bigl(-c'(\log x)^{3/5}(\log\log x)^{-1/5}\bigr).
\end{equation}
\end{corollary}

\begin{proof}
Write $\phi(y)=(\log y)^{3/5}(\log\log y)^{-1/5}$.  On an excursion ending
at $v_\nu$, \eqref{eq:VK_M} gives
$H_\nu\ll v_\nu e^{-c\phi(v_\nu)}$, hence
$\sqrt{H_\nu}\ll v_\nu^{1/2}e^{-c\phi(v_\nu)/2}$.  The same pointwise
majorant applies to the corresponding triangular skeleton, with a possibly
smaller exponential constant because $\phi$ is eventually increasing and
slowly varying on fixed multiplicative intervals.  Integrating the resulting
majorant up to $x$ gives \eqref{eq:uncond_area} for some $c'>0$; no identity
$c'=c/2$ is asserted.
\end{proof}

\begin{corollary}
\label{cor:rh}
\emph{[CONDITIONAL: RH]} If RH holds, then
\begin{equation}
\label{eq:rh_area}
  \Ap(x)+|\Am(x)|\ll x^{5/4+\eps}.
\end{equation}
Consequently, if there exists a fixed $\delta>0$ such that
$\Ap(x)+|\Am(x)|\gg x^{5/4+\delta}$ along an unbounded sequence, RH is
false.
\end{corollary}

\begin{remark}[False inference withdrawn]
\label{rem:onedirection}
\textbf{Status: FALSE / WITHDRAWN.} From $\ell_\nu\ge2H_\nu$ and a bound
$\sqrt{H_\nu}\ell_\nu\ll v_\nu^{5/4+\eps}$ one may deduce only
$H_\nu\ll v_\nu^{5/6+O(\eps)}$. This is a bound indexed by the
\emph{future endpoint} $v_\nu$; it does \emph{not} imply
$M(x)\ll x^{5/6+\eps}$ at interior points, and therefore gives no
zero-free half-plane. Version 5's claimed consequence
$\operatorname{Re}s>5/6$ was invalid. Whether \eqref{eq:rh_area} is
equivalent to RH remains unresolved by the present argument.
\end{remark}

\begin{remark}[Oscillation of the heights]
\label{rem:omega}
The strongest published numerical lower bounds recorded here are
$\limsup M(x)/\sqrt{x}>1.826054$ and
$\liminf M(x)/\sqrt{x}<-1.837625$ \cite{Hurst18}. In particular there are
infinitely many $n$ with $M(n)>1.826054\sqrt n$; the excursion containing $n$
has $u_\nu<n$, so $H_\nu>1.826054\,u_\nu^{1/2}$ for infinitely many
positive excursions, and similarly with $1.837625$ for negative ones. Version 5 claimed
$H_\nu\gg v_\nu^{1/2-\eps}$; this is \emph{not justified}, since $v_\nu$
may be much larger than $n$.
\end{remark}

\section{Criteria equivalent to RH}
\label{sec:criteria}

\noindent\textbf{Status: MIXED.} The Mellin identity and the global
$L^1/L^2$ equivalences are classical consequences of Littlewood's criterion
and partial summation; they are \emph{not claimed as new RH criteria}. The
excursion decompositions, pruning lemmas and local sufficient conditions are
exact elementary repackagings proved here. Heuristics are explicitly
unproved.

\medskip

For $\operatorname{Re}s>1$,
\begin{equation}
\label{eq:mellin_M}
  \int_1^\infty M(x)x^{-s-1}\,dx=\frac{1}{s\,\zeta(s)} .
\end{equation}
If $\int_1^\infty|M(x)|x^{-\sigma-1}dx<\infty$ for every $\sigma>1/2$, the
left side continues analytically to $\operatorname{Re}s>1/2$, so $\zeta$ has
no zeros there.

\begin{proposition}[$L^1$ criterion]
\label{prop:L1}
\emph{[PROVED HERE]} RH holds if and only if, for every $\eps>0$,
\begin{equation}
\label{eq:L1_criterion}
  \sum_{v_\nu\le x}\beta_\nu+\beta_{\rm term}(x)=\sum_{n\le x}|M(n)|
  \ll x^{3/2+\eps}.
\end{equation}
\end{proposition}

\begin{proof}
RH gives $M(x)\ll_\eps x^{1/2+\eps}$
\cite[\S14.25]{Titchmarsh1986}. Conversely, fix any $\sigma>1/2$ and choose
$0<\eps<\sigma-1/2$.  Partial summation applied to
\eqref{eq:L1_criterion} then gives
$\sum_n|M(n)|n^{-\sigma-1}<\infty$; since $M$ is constant on $[n,n+1)$,
apply the remark after \eqref{eq:mellin_M}.
\end{proof}

\begin{lemma}[Pruning]
\label{lem:pruning}
\emph{[PROVED HERE]} $\beta_\nu\le H_\nu\ell_\nu\le\ell_\nu^2/2$, hence
\begin{equation}
\label{eq:pruning}
  \sum_{v_\nu\le x,\ \ell_\nu\le L_0}\beta_\nu\le\frac{xL_0}{2}.
\end{equation}
\end{lemma}

\begin{corollary}[Reduction to long excursions]
\label{cor:long}
\emph{[PROVED HERE]} For fixed $\delta\ge0$ and $L_0=x^{1/2-\delta}$, RH
holds if and only if
\begin{equation}
\label{eq:long_reduction}
  \sum_{\substack{v_\nu\le x\\ \ell_\nu>L_0}}\beta_\nu
  +\mathbf 1_{\{L_x>L_0\}}\beta_{\rm term}(x)\ll x^{3/2+\eps}
  \quad\text{for every }\eps>0.
\end{equation}
The omitted terminal block when $L_x\le L_0$ is bounded by the same trivial
estimate as the omitted complete excursions. At most
$x/L_0$ excursions are involved.
\end{corollary}

\begin{remark}[The $\sqrt{x}$ frontier]
\label{rem:sqrt_frontier}
The cut may even depend on $\eps$: with $L_0=x^{1/2+\eps/2}$ the discarded
part is $\le x^{3/2+\eps/2}/2$. The frontier is therefore $x^{1/2+o(1)}$;
what requires new information is a cut $x^{\theta}$ with a fixed
$\theta>1/2$ (version 5 stated the frontier as $x^{1/2}$).
\end{remark}

\begin{heuristic}[Renewal tail]
\label{heu:renewal}
This is a heuristic argument. If excursion lengths behave like those of a
recurrent random walk,
\begin{equation}
\label{eq:renewal_tail}
  \#\{v_\nu\le x:\ \ell_\nu>L\}\asymp\sqrt{x}\,L^{-1/2},
\end{equation}
and with $\beta_\nu\asymp c_\beta\ell_\nu^{3/2}$ the dyadic range
$[L,2L]$ contributes $\asymp\sqrt{x}\,L$ to \eqref{eq:L1_criterion}, so the
terminal scales dominate.
\end{heuristic}

\begin{remark}[Limit of pruning]
\label{rem:pruning_limit}
For $\ell_\nu\asymp x$ the trivial bound gives $\beta_\nu\ll x^2$ against a
budget $x^{3/2+\eps}$: a factor $x^{1/2}$ is missing, and there are only
$O(1)$ such excursions.
\end{remark}

\begin{theorem}[Correlation criterion]
\label{thm:corr-criterion}
\emph{[PROVED HERE]} RH holds if and only if, for every $\eps>0$,
\begin{equation}
\label{eq:RH-corr}
  \sum_{n\le x}M(n)^2=\sum_{v_\nu\le x}Q_\nu+Q_{\rm term}(x)
  =-\sum_{v_\nu\le x}\sum_{h}hC^{(\nu)}_h+Q_{\rm term}(x)\ll x^{2+\eps}.
\end{equation}
\end{theorem}

\begin{proof}
Under RH apply Littlewood with exponent $\eps/2$, so
$M(n)^2\ll_\eps n^{1+\eps}$ and the required quadratic sum is
$O_\eps(x^{2+\eps})$. Conversely Cauchy--Schwarz gives
$\sum_{n\le x}|M(n)|\le x^{1/2}(\sum M^2)^{1/2}\ll x^{3/2+\eps/2}$; apply
Proposition~\ref{prop:L1}.
\end{proof}

\begin{corollary}[Local sufficient condition, version 5]
\label{cor:local-corr-sufficient}
\emph{[PROVED HERE: unconditional implication]} If for every $\eps>0$ there is
$C_\eps$ with
\begin{equation}
\label{eq:local_Q}
  Q_\nu\le C_\eps\,\ell_\nu^{2+\eps}\quad\text{for every complete excursion},
\end{equation}
and, separately, every terminal block satisfies
\[
Q_{\rm term}(x)\le C_\eps L_x^{2+\eps},
\]
then RH holds.
\end{corollary}

\begin{heuristic}[The hypothesis of Corollary~\ref{cor:local-corr-sufficient} fails]
\label{heu:local_fails}
This is a heuristic argument. In a model where, after a zero of $M$, the next
values of $\mu$ behave locally at random with squarefree density $6/\pi^2$,
the pattern ``$k$ non-zero steps of one sign followed by $k$ of the other''
follows a given zero with probability $\ge c_0^{-k}$. Even with infinitely many zeros, this model does \emph{not} by itself
prove that such patterns occur for unbounded $k$; that would require a
quantitative number of trials together with dependence control. The
statement that excursions with $H_\nu=k$ and $\ell_\nu\asymp k$ should
occur for $k\to\infty$ is therefore only a Borel--Cantelli-style
heuristic prediction. By \eqref{eq:height_Q_min}
$Q_\nu/\ell_\nu^{2+\eps}\gg k^{1-\eps}\to\infty$, so \eqref{eq:local_Q}
fails for every $\eps\in(0,1)$.  For a fixed $\eps\ge1$ the individual inequality
is trivial from \eqref{eq:Q_upper}; this does \emph{not} make the hypothesis
``for every $\eps>0$'' true. Version 5 called the condition ``sufficient but
not weak''; heuristically it is never satisfied.
\end{heuristic}

\begin{corollary}[Pruned local condition]
\label{cor:pruned_local}
\emph{[PROVED HERE: unconditional implication]} If for every $\eps>0$ there is $C_\eps$
with
\begin{equation}
\label{eq:pruned_Q}
  Q_\nu\le C_\eps\,\ell_\nu^{2+\eps}\quad\text{whenever }\ell_\nu>v_\nu^{1/2}
\end{equation}
(and likewise for the terminal block when its length exceeds $x^{1/2}$),
then RH holds.
\end{corollary}

\begin{proof}
Split the completed excursions at $\ell_\nu=x^{1/2}$.  If
$\ell_\nu\le x^{1/2}$, then \eqref{eq:Q_upper} gives
$Q_\nu\le(\ell_\nu^3+2\ell_\nu)/12\le\ell_\nu x/12+\ell_\nu/6$; disjointness
of the interiors makes their total contribution at most $x^2/12+x/6$.
If $\ell_\nu>x^{1/2}$, then automatically
$\ell_\nu>v_\nu^{1/2}$ because $v_\nu\le x$, so the stated hypothesis
applies and gives at most $C_\eps\sum\ell_\nu^{2+\eps}\le C_\eps x^{2+\eps}$.
Treat the terminal block by the separately stated terminal hypothesis and
apply Theorem~\ref{thm:corr-criterion}.
\end{proof}

\noindent The ballistic excursions of Heuristic~\ref{heu:local_fails} have
$\ell_\nu\asymp k\ll v_\nu^{1/2}$ and do not contradict \eqref{eq:pruned_Q}.

\begin{remark}[Scale of the correlation sums]
\label{rem:chowla-scale}
Under the crude model in which ``$j$ squarefree'' and ``$j+h$ squarefree''
are independent with probability $6/\pi^2$ each,
$P_h\approx(6/\pi^2)^2(L-h)$ and
\begin{equation}
\label{eq:chowla_crude}
  \sum_{h<L}hP_h\approx\Bigl(\frac{6}{\pi^2}\Bigr)^2\sum_{h<L}h(L-h)
  \approx\frac{6}{\pi^4}L^3\approx0.0616\,L^3 .
\end{equation}
Version 5 gave $\tfrac12(6/\pi^2)^2L^3\approx0.18L^3$, omitting the factor
$L-h$; that constant is \emph{incorrect}, the order $L^3$ stands. For
$|\sum_hhC_h|$ of size $B(\sum_hh^2)^{1/2}\asymp BL^{3/2}$ to be $O(L^2)$ one
needs $B=O(\sqrt L)$ (a calibration, not a proof).
\end{remark}

\begin{remark}[Conditioning on the bridge]
\label{rem:conditioning}
Inside an excursion the correlations are constrained by
\eqref{eq:Qcorr}, \eqref{eq:zeroth-corr} and by the geometry of the bridge;
with the defect $d_\nu:=\ell_\nu/2-H_\nu\ge0$, a third constraint of this
type is a research direction, not a result.
\end{remark}

\begin{theorem}[Energy criterion]
\label{thm:energy}
\emph{[PROVED HERE]} Define, with $u$ the last zero $\le x$ and
$H_{\rm run}(x)=\max_{u<n\le x}|M(n)|$,
\begin{equation}
\label{eq:energy_def}
  \En(x):=\frac14\Bigl(\sum_{v_\nu\le x}H_\nu\ell_\nu+(x-u)H_{\rm run}(x)\Bigr).
\end{equation}
For a complete excursion let $J$ be the chord joining
$(u_\nu,\Anet(u_\nu))$ to $(v_\nu,\Anet(v_\nu))$.  By
\eqref{eq:jump}, $J'=\sigma_\nu\sqrt{H_\nu}/2$, and therefore
$\tfrac14H_\nu\ell_\nu=\int_{I_\nu}(J')^2$; for the
terminal block this interpretation does not hold and \eqref{eq:energy_def} is
a definition. Then
\begin{equation}
\label{eq:energy_sandwich}
  \sum_{n\le x}|M(n)|\le4\En(x)\le x\max_{n\le x}|M(n)|,
\end{equation}
and RH holds if and only if $\En(x)\ll x^{3/2+\eps}$ for every $\eps>0$.
\end{theorem}

\begin{proof}
$\beta_\nu\le H_\nu\ell_\nu$ and $\sum\ell_\nu\le x$; apply
Proposition~\ref{prop:L1}.
\end{proof}


\section{Lacunary long excursions and localized RH criteria}
\label{sec:lacunary}

\noindent\textbf{Status: [PROVED HERE] with [PUBLISHED INPUT] from Kaczorowski--Pintz.} Lemmas~\ref{lem:power-lacunary} and
\ref{lem:linear-lacunary} are unconditional geometric counting bounds,
using only disjointness and ordering of excursions. The arithmetic input is
Corollary~4 of Kaczorowski--Pintz~\cite{KP86}: at $a=0$ their function is
$M_0=M+2$, not $M$. The proof below explicitly accounts for this additive
constant. Subject to that published theorem, the localized RH criteria below
are unconditional theorems.

Put
\[
\Theta=\sup\{\Re\rho:\zeta(\rho)=0,\ 0<\Re\rho<1\},\qquad
\gamma_1=14.134725\ldots .
\]

\begin{definition}[Power- and linear-long excursions]
For $A>1$, $I_\nu=(x_{0,\nu},x_{1,\nu}]$ is \emph{$A$-exceptional} if
$x_{1,\nu}>x_{0,\nu}^A$ and $A$-ordinary otherwise. For $\lambda>1$ it is
\emph{$\lambda$-long} if $x_{1,\nu}>\lambda x_{0,\nu}$. The initial
excursion $(0,2]$ is treated as one finite exceptional case.
\end{definition}

\begin{lemma}[Power lacunarity: the $O(\log\log X)$ bound]
\label{lem:power-lacunary}
For every $A>1$,
\begin{equation}
\label{eq:power-lacunary-count}
E_A(X):=\#\{\nu:x_{0,\nu}\le X,\ x_{1,\nu}>x_{0,\nu}^A\}
\le2+\frac1{\log A}\log\frac{\log X}{\log2}.
\end{equation}
Thus $E_A(X)=O_A(\log\log X)$. In particular, for $A=2$,
\begin{equation}
\label{eq:square-exception-count}
E_2^{\rm pow}(X)\le2+\log_2\frac{\log X}{\log2}.
\end{equation}
\end{lemma}

\begin{proof}
List exceptional starts at least $2$ as $2\le a_1<a_2<\cdots$. Since the
next excursion starts at or after the end of the preceding one, possibly
after a zero run,
the end $x_{1,j}$ of the preceding exceptional excursion satisfies
\[
a_{j+1}\ge x_{1,j}>a_j^A.
\]
Hence $\log a_j>A^{j-1}\log2$.
\end{proof}

\begin{lemma}[Linear lacunarity: the $O(\log X)$ bound]
\label{lem:linear-lacunary}
For every $\lambda>1$ and $X\ge2$,
\begin{equation}
\label{eq:linear-lacunary-count}
E_\lambda^{\rm lin}(X):=
\#\{\nu:x_{0,\nu}\le X,\ x_{1,\nu}>\lambda x_{0,\nu}\}
\le2+\frac{\log(X/2)}{\log\lambda}.
\end{equation}
In particular the excursions with $x_1>2x_0$ are at most
$(\log2)^{-1}\log X+O(1)$.
\end{lemma}

\begin{proof}
For starts $a_j\ge2$, disjointness gives $a_{j+1}>\lambda a_j$, hence
$a_j>2\lambda^{j-1}$.
\end{proof}

\begin{theorem}[Kaczorowski--Pintz, specialized to $M$]
\label{thm:KP-M}
\textbf{Status: [PUBLISHED INPUT].} Let $M_0(x)=M(x)+2$.
Kaczorowski--Pintz~\cite[Corollary~4, p.~178, equations (2.13)--(2.15)]{KP86}
define $V(f,Y)$ by alternating-sign points and state
\begin{equation}
\label{eq:KP-signchanges}
\liminf_{Y\to\infty}\frac{V(M_0,Y)}{\log Y}
\ge\frac{\gamma_1}{\pi},
\end{equation}
with ``combined oscillation of size'' $x^{\Theta-\eta}$ for every $\eta>0$,
in the sense of their definition following (1.1): the selected alternating
points $t_j$ additionally satisfy $|M_0(t_j)|\ge t_j^{\Theta-\eta}$.
Moreover, Corollary~4 is obtained from their Corollary~1, and
Corollary~1 states that relations (2.1)--(2.3) of Theorem~1 hold. Hence, for
every $\delta>0$, every interval
\[
[Y^{1-\delta},Y]\qquad(Y>Y_0(\delta))
\]
contains a sign change of $M_0$ combined with the same oscillation size
(Theorem~1, equations (2.2)--(2.3), pp.~175--176).
Consequently, if $\Theta>1/2$ and $0<\eta<\Theta-1/2$, then
\begin{equation}
\label{eq:KP-large-excursions}
N_\eta(Y)\ge\left(\frac{\gamma_1}{\pi}-o(1)\right)\log Y,
\end{equation}
where $N_\eta(Y)$ counts distinct excursions starting at most $Y$ that
contain a point $t\le Y$ with
\[
|M(t)|\ge\tfrac12 t^{\Theta-\eta}.
\]
\end{theorem}

\begin{proof}
At $a=0$, the function $M_a$ of~\cite{KP86} is
$M_0=M-1/\zeta(0)=M+2$. The combined oscillation points alternate sign and
have $|M_0(t)|\ge t^{\Theta-\eta}$. For sufficiently large $t$ the positive
points have $M(t)>0$ and the negative points have $M(t)<0$; therefore
consecutive selected points, and hence all selected alternating points, lie
in distinct excursions. Their starts are below the points and thus at most
$Y$. For all sufficiently large selected points,
\[
|M(t)|\ge |M_0(t)|-2\ge \tfrac12 t^{\Theta-\eta},
\]
so the additive constant $2$ is absorbed without changing the exponent.
\end{proof}

\begin{corollary}[Published multiplicative localization of Mertens returns]
\label{cor:KP-local-returns}
\textbf{Status: [PROVED HERE FROM PUBLISHED INPUT].} For every fixed $A>1$
there exists $U_A$ such that every excursion starting at $u\ge U_A$ ends at
$v\le u^A$. Equivalently, the actual Mertens function has only finitely many
$A$-exceptional excursions.
\end{corollary}

\begin{proof}
Fix $A>1$ and put $\delta=1-1/A$. Apply the local part of
Theorem~\ref{thm:KP-M} with $Y=u^A$; its interval is
$[Y^{1-\delta},Y]=[u,u^A]$. Choose $\eta_0$ with
$0<\eta_0<\Theta$; this is possible because $\Theta\ge1/2$. For $u$
sufficiently large, the two alternating-sign witnesses for the local sign
change of $M_0=M+2$ have magnitude $>2$. Thus the positive witness has
$M>0$ and the negative witness has $M<0$. Since the integer-valued Mertens
function changes only by $0,\pm1$, it takes the value $0$ at an integer
strictly between those witnesses. If $u=x_{0,\nu}$ is the beginning of an
excursion, its first return $x_{1,\nu}$ is no later than that zero, and
therefore $x_{1,\nu}\le u^A$.
\end{proof}

\begin{remark}[Geometric count versus published localization]
Lemma~\ref{lem:power-lacunary} is an elementary statement about any ordered
family of disjoint intervals and gives the independent estimate
$E_A(X)=O_A(\log\log X)$. Corollary~\ref{cor:KP-local-returns} is stronger
for the actual Mertens excursions: for each fixed $A>1$, the exceptional
family is finite. The geometric estimate remains a separate theorem because
it requires no arithmetic input.
\end{remark}

\begin{theorem}[Sparse-exception localization of RH]
\label{thm:sparse-RH}
\textbf{Status: [PROVED HERE, USING PUBLISHED INPUT].} Let $\mathcal E$ be any family of excursions satisfying
\[
\#\{\nu\in\mathcal E:x_{0,\nu}\le X\}=o(\log X).
\]
Then RH is equivalent to: for every $\varepsilon>0$ there exists
$C_\varepsilon$ such that
\begin{equation}
\label{eq:restricted-RH}
|M(t)|\le C_\varepsilon t^{1/2+\varepsilon}
\end{equation}
for every $t$ lying in every excursion outside $\mathcal E$. No bound of any
kind is required on the excursions in $\mathcal E$.
\end{theorem}

\begin{proof}
RH implies the global estimate by Littlewood. Conversely, suppose
$\Theta>1/2$ and choose $\eta,\varepsilon>0$ with
$1/2+\varepsilon<\Theta-\eta$. Theorem~\ref{thm:KP-M} supplies
$\gg\log X$ distinct excursions containing values
$\gg t^{\Theta-\eta}$, whereas only $o(\log X)$ may belong to $\mathcal E$.
More explicitly,
\[
\left(\frac{\gamma_1}{\pi}-o(1)\right)\log X-o(\log X)\longrightarrow\infty,
\]
so arbitrarily large selected excursions lie outside $\mathcal E$.  At their
large points,
\[
\frac{t^{\Theta-\eta}}{t^{1/2+\varepsilon}}
=t^{\Theta-\eta-1/2-\varepsilon}\longrightarrow\infty,
\]
contradicting \eqref{eq:restricted-RH}.
\end{proof}

\begin{corollary}[Power-localized RH criterion]
\label{cor:power-ordinary-RH}
\textbf{Status: [PROVED HERE].} For every fixed $A>1$, RH is equivalent to \eqref{eq:restricted-RH}
restricted only to $A$-ordinary excursions $x_1\le x_0^A$. In particular,
\begin{equation}
\boxed{\mathrm{RH}\iff
|M(t)|\ll_\varepsilon t^{1/2+\varepsilon}
\ \text{for all $t$ in excursions with }x_1\le x_0^2.}
\label{eq:square-horizon-RH}
\end{equation}
No estimate at all is required on the excursions with $x_1>x_0^2$.
\end{corollary}

\begin{proof}
Lemma~\ref{lem:power-lacunary} gives $E_A(X)=O_A(\log\log X)=o(\log X)$.
Apply Theorem~\ref{thm:sparse-RH}.
\end{proof}

\begin{corollary}[Square-truncated excursion criterion]
\label{cor:square-truncated-RH}
\textbf{Status: [PROVED HERE].} RH is equivalent to the following statement:
for every $\varepsilon>0$ there exists $C_\varepsilon$ such that, for every
excursion $I_\nu=(x_{0,\nu},x_{1,\nu}]$,
\[
|M(t)|\le C_\varepsilon t^{1/2+\varepsilon}
\qquad
\left(x_{0,\nu}<t\le\min\{x_{1,\nu},x_{0,\nu}^2\}\right).
\]
\end{corollary}

\begin{proof}
For a square-ordinary excursion, $x_{1,\nu}\le x_{0,\nu}^2$, so the
truncated range is the whole excursion.  Apply
Corollary~\ref{cor:power-ordinary-RH}.  On square-exceptional excursions the
criterion deliberately imposes no condition beyond $x_0^2$.
\end{proof}

\begin{remark}[What the $O(\log\log X)$ theorem does and does not bound]
\label{rem:no-height-from-lacunarity}
\textbf{Status: [NO USEFUL BOUND] on individual height from lacunarity alone.}
The theorem $E_A(X)=O_A(\log\log X)$ bounds only the \emph{number} of
$A$-exceptional excursions. It gives no upper or lower RH-scale bound for
$H_\nu$, $\beta_\nu$ or $Q_\nu$ of an individual exceptional excursion.
No probabilistic tail estimate is used anywhere in
Corollary~\ref{cor:power-ordinary-RH}. In particular the previously
mentioned $\sqrt{\log\log\log X}$ extreme-value scale is a heuristic
from a random model, not a theorem, and is not used in this manuscript.
\end{remark}

\begin{corollary}[Linear-localized RH criterion]
\label{cor:linear-ordinary-RH}
\textbf{Status: [PROVED HERE].} Let
\[
\lambda_*:=\exp(\pi/\gamma_1).
\]
For every fixed $\lambda>\lambda_*$, RH is equivalent to
\eqref{eq:restricted-RH} restricted only to excursions with
$x_1\le\lambda x_0$. In particular it is enough to control excursions with
$x_1\le2x_0$.
\end{corollary}

\begin{proof}
The complementary family has at most
$(1/\log\lambda+o(1))\log X$ members by
Lemma~\ref{lem:linear-lacunary}. If RH fails, Theorem~\ref{thm:KP-M} gives
at least $(\gamma_1/\pi-o(1))\log X$ large distinct excursions. The strict
inequality $1/\log\lambda<\gamma_1/\pi$ leaves a positive proportion outside
the long family.
\end{proof}

\begin{remark}[Role of the two new counts]
\label{rem:two-counts}
The $O(\log\log X)$ and $O(\log X)$ bounds are geometric, not probabilistic.
Their arithmetic force appears after combination with the independent
Kaczorowski--Pintz lower bound of order $\log X$ for large alternating
oscillations. The exponent $2$ is not needed for localization itself; it is
singled out next because the range $t\le x_0^2$ has an exact bilinear kernel.
\end{remark}

\subsection{The square range as an exact bilinear form of the prefix}

\begin{proposition}[Bilinear square-range identity]
\label{prop:prefix-bilinear}
Let $u\ge1$ and $1\le t\le u^2$. Then
\begin{equation}
\label{eq:bilinear-general}
M(t)=2M(\min(t,u))-
\sum_{d,j\le u}\mu(d)\mu(j)\left\lfloor\frac{t}{dj}\right\rfloor.
\end{equation}
In particular, if $M(u)=0$ and $u<t\le u^2$,
\begin{equation}
\label{eq:bilinear-excursion}
M(t)=-\mathcal B_u(t),\qquad
\mathcal B_u(t):=\sum_{d,j\le u}\mu(d)\mu(j)
\left\lfloor\frac{t}{dj}\right\rfloor,
\qquad \mathcal B_u(u)=0.
\end{equation}
\end{proposition}

\begin{proof}
Let $\mu_u=\mu\mathbf1_{[1,u]}$. The convolution $(\mu-\mu_u)*(\mu-\mu_u)$ is supported strictly above
$u^2$ because each factor is supported on integers $>u$; convolution with
$1$ cannot move support to a smaller integer.  Hence
$(\mu-\mu_u)*(\mu-\mu_u)*1$ vanishes for $n\le u^2$. Expanding and
using $\mu*1=\delta$ gives
$\mu(n)=2\mu_u(n)-(\mu_u*\mu_u*1)(n)$ for $n\le u^2$; summing gives the
formula.
\end{proof}

\begin{remark}[Higher-order determinacy; why the square is special]
\label{rem:prefix-k}
For every $k\ge1$ and $n\le u^k$,
\begin{equation}
\label{eq:HB-k}
\mu(n)=\sum_{j=1}^{k}(-1)^{j-1}\binom{k}{j}
\bigl(\mu_u^{*j}*1^{*(j-1)}\bigr)(n).
\end{equation}
Indeed, $(\mu-\mu_u)^{*k}$ is supported strictly above $u^k$, and therefore
\[
(\mu-\mu_u)^{*k}*1^{*(k-1)}(n)=0\qquad(n\le u^k).
\]
Expand the left side by the binomial theorem and repeatedly use
$\mu*1=\delta$; the $j=0$ term reduces to $\mu$, and rearrangement gives
\eqref{eq:HB-k}. Thus ``determined by the prefix'' is not special to $u^2$. What is special
to $t\le u^2$ is the order-$2$ form \eqref{eq:bilinear-excursion}, with
rank-one coefficients $\mu(d)\mu(j)$ and elementary kernel
$\lfloor t/(dj)\rfloor$. This is a Heath--Brown/Vaughan-type truncation identity; related identities
and summed forms occur in the cited literature and in fast computation of
$M$ \cite{HB82,HT23}.  No claim is made here that the exact displayed
normalization appears verbatim in those references.
\end{remark}

\begin{question}[Central square-excursion problem]
\label{q:excursion-constraint}
\textbf{Status: OPEN.}\par\noindent\textbf{No RH-scale bound is proved here.}
Let $u=x_{0,\nu}$ and $x_{1,\nu}\le u^2$. Then
\[
\mathcal B_u(u)=\mathcal B_u(x_{1,\nu})=0,
\]
and $-\sigma_\nu \mathcal B_u(t)>0$ for
$u<t<x_{1,\nu}$. Can the rank-one coefficient matrix $\mu(d)\mu(j)$, the
product kernel $\lfloor t/(dj)\rfloor$, and these bridge/sign constraints be
combined to prove
\[
\max_{u<t\le x_{1,\nu}}|\mathcal B_u(t)|
\ll_\varepsilon x_{1,\nu}^{1/2+\varepsilon}
\]
uniformly over all such excursions? The uniform mathematical bound stated here is equivalent to RH by
Corollary~\ref{cor:power-ordinary-RH}.  The additional phrase ``prove it from
the bilinear structure'' is a research program, not another equivalent
mathematical statement.  The formulation isolates that program inside
finite bilinear bridges of the square range.
\end{question}

\subsection{Mellin mass of a complete excursion: the direction $M\to\zeta$}

\begin{proposition}[Exact excursion block and its Mellin mass]
\label{prop:excursion-mellin}
For a complete excursion $I=(x_0,x_1]$, put
\[
D_I(s)=\sum_{x_0<n\le x_1}\frac{\mu(n)}{n^s}.
\]
Then
\begin{equation}
\label{eq:excursion-abel}
D_I(s)=\sum_{n=x_0+1}^{x_1-1}M(n)
\bigl(n^{-s}-(n+1)^{-s}\bigr).
\end{equation}
For real $\sigma>0$, $D_I(\sigma)$ has the sign of the excursion and
\begin{equation}
\label{eq:excursion-mass-lower}
|D_I(\sigma)|\ge
\sigma\beta_I x_1^{-\sigma-1}
\ge\sigma H_I^2x_1^{-\sigma-1}.
\end{equation}
For $\Re s>1$,
\begin{equation}
\label{eq:1overzeta-excursions}
\frac1{\zeta(s)}=\sum_I D_I(s),
\end{equation}
zero runs contributing nothing.
\end{proposition}

\begin{proof}
Summation by parts gives \eqref{eq:excursion-abel} because the endpoint
values of $M$ vanish. On an excursion all interior $M(n)$ have the same
sign, and
$n^{-\sigma}-(n+1)^{-\sigma}\ge\sigma x_1^{-\sigma-1}$. The bound
$\beta_I\ge H_I^2$ follows from Proposition~\ref{prop:sharp-LH}. If $M$ is zero on two consecutive integer points, its intervening increment
$\mu(n)$ is zero; hence zero runs contribute no Dirichlet terms.  Therefore
all integers with $\mu(n)\ne0$ belong to the excursion blocks, and summing
the disjoint blocks gives \eqref{eq:1overzeta-excursions} in the half-plane
of absolute convergence.
\end{proof}

\begin{remark}[No finite-data assumption]
\label{rem:no-finite-data}
Nothing in this section treats the computed range of $M$ as representative.
The statements are uniform in $x_0$; a dangerous excursion may begin at
$10^{40}$ or arbitrarily farther out. The role of $x_0^2$ is algebraic, not
empirical.
\end{remark}

\section{Trial-division horizon of an excursion}
\label{sec:prime_horizon}

\noindent\emph{Status.} Lemma~\ref{lem:square-horizon} is standard trial
division. The $O(\log\log X)$ separation has already been isolated in
Lemma~\ref{lem:power-lacunary}; Proposition~\ref{prop:double-exp} records its
$A=2$ form in the historical notation of versions 5--6.

\begin{lemma}[Square horizon]
\label{lem:square-horizon}
\emph{[PROVED HERE]} Let $u\ge2$. Trial division by the primes $p\le u$
alone determines $\mu(n)$ for every $n\le u^2$.
\end{lemma}

\begin{proof}
Remove the prime factors $p\le u$ of $n$; if one occurs twice, $\mu(n)=0$.
The cofactor $q\le n\le u^2$ has all prime factors $>u$, so it is $1$ or a
single prime (two such factors would give $q>u^2$).
\end{proof}

\begin{warning}[On the word ``information'']
\label{warn:information_language}
Primes $>u$ are not logically new information, and prefix determinacy is not
special to $u^2$: the higher-order identity \eqref{eq:HB-k} propagates a
finite prefix to every fixed power range.  The square is special only in the
more concrete sense that one-stage trial division by primes $\le u$ suffices
and the exact summatory formula has the bilinear kernel
$\lfloor t/(dj)\rfloor$.  Historical uses of ``known'', ``new'' and
``information'' below refer only to this finite-level decomposition, not to
logical independence.
\end{warning}

\begin{definition}[Smooth and rough parts]
\label{def:smooth_rough}
For $y>1$ write $n=s_y(n)\,r_y(n)$, where every prime factor of $s_y(n)$
is $\le y$ and every prime factor of $r_y(n)$ is $>y$. Let $\Pp(n)$ be the
largest prime factor of $n$, with $\Pp(1)=1$.
\end{definition}

\begin{lemma}
\label{lem:rough_part}
\emph{[PROVED HERE]}
\begin{equation}
\label{eq:mu_split}
  \mu(n)=\mu\bigl(s_y(n)\bigr)\,\mu\bigl(r_y(n)\bigr);
\end{equation}
if $k\in\mathbb N$ and $n<y^k$ then $\Omega(r_y(n))\le k-1$, and if $n\le y^2$ then $r_y(n)$ is
$1$ or a prime.
\end{lemma}

\begin{proof}
The two parts are coprime. If $\Omega(r)\ge1$ then $y^{\Omega(r)}<r\le n$.
\end{proof}

\begin{definition}[Square-horizon classes]
\label{def:info-excursion}
An excursion is \emph{square-ordinary} if $v_\nu\le u_\nu^2$ and
\emph{square-exceptional} otherwise. These replace the historical information terminology of versions 5--6.\par
For $u_\nu\ge2$,
\[
\vartheta_\nu=\min\left\{1,\frac{u_\nu^2-u_\nu}{\ell_\nu}\right\}
\]
is the fraction of the excursion lying below the square horizon.
\end{definition}

\begin{proposition}[Double-exponential separation]
\label{prop:double-exp}
Let $a_1<a_2<\cdots$ be the starting points at least $2$ of
square-exceptional excursions. Then
\[
a_{j+1}>a_j^2,\qquad a_j\ge a_1^{2^{j-1}},
\]
and
\[
N_{\rm sq}(X)=\#\{j:a_j\le X\}
\le1+\log_2\frac{\log X}{\log2}.
\]
This is the $A=2$ case of Lemma~\ref{lem:power-lacunary}.
\end{proposition}

\begin{proof}
The end of a square-exceptional excursion starting at $a_j$ is beyond
$a_j^2$, and the next such excursion starts at or after that end.
\end{proof}

\begin{corollary}[Square-exceptional excursions are finite]
\label{cor:finiteness_expanding}
\textbf{Status: [PROVED HERE FROM PUBLISHED INPUT].} There are only finitely
many square-exceptional excursions. More generally, for every fixed $A>1$
there are only finitely many excursions with $x_1>x_0^A$.
\end{corollary}

\begin{proof}
Take $A=2$ in Corollary~\ref{cor:KP-local-returns}.
\end{proof}

\begin{remark}
A bare sign change of $M_0=M+2$ would not force a zero of $M$. The decisive
extra point is the published phrase ``combined oscillation of size'': the
local sign change comes with two opposite-sign witnesses whose magnitudes
tend to infinity. Once they exceed $2$, the corresponding values of $M$
have opposite signs, and the unit-step property forces an intervening zero.
\end{remark}

\begin{remark}[The horizon $u^2$ is not the multiplicative threshold of an excursion]
\label{rem:horizon_inactive}
For a square-ordinary excursion $\sqrt{v_\nu}\le u_\nu$. Every $n\le v_\nu$
has at most one prime factor $>\sqrt{v_\nu}$, so the threshold that governs
the multiplicative structure inside $I_\nu$ is $\sqrt{v_\nu}$, which lies
below $u_\nu$. The split ``known/new'' at $u_\nu$ does not match it, and for
a square-ordinary excursion the regime $t>u_\nu^2$ of
Section~\ref{sec:prime_copy_propagation} lies beyond $v_\nu$. That section
is therefore about $M(t)$, not about a single excursion;
Section~\ref{sec:local_large_prime} gives the local version.
\end{remark}

\section{Large-prime decomposition at level \texorpdfstring{$u^2$}{u^2}}
\label{sec:prime_copy_propagation}

\noindent\textbf{Status: [PROVED HERE], with classical sieve ancestry.} The
identities are elementary decompositions by large prime factors and are
instances of Buchstab/Legendre-type sieve organization; no blanket claim of
originality is made. The asymptotic statements use only the prime number
theorem with classical error term and the Vinogradov--Korobov bound
\cite{MontgomeryVaughan2007,Walfisz1963}, and are proved here. Several
statements of version 5 are corrected; see
Warnings~\ref{warn:v5_Ny}, \ref{warn:v5_main_term},
\ref{warn:margin_vacuous} and~\ref{warn:v5_calibration}.

\begin{definition}[Smooth baseline and large-prime part]
\label{def:baseline_channel}
For $u\ge2$ and $t\ge1$ put
\begin{equation}
\label{eq:baseline_channel_def}
  S_u(t):=\sum_{\substack{n\le t\\ \Pp(n)\le u^2}}\mu(n),\qquad
  E_u(t):=\sum_{\substack{n\le t\\ \Pp(n)>u^2}}\mu(n),\qquad
  M(t)=S_u(t)+E_u(t).
\end{equation}
Version 5 used an undefined ``baseline'' $S_u(t)$; \eqref{eq:baseline_channel_def}
is the definition adopted here. In the window $u^2<t<u^4$, $E_u$ is the
quantity $E_u^{(1)}$ of version 5.
\end{definition}

\begin{proposition}[Single-prime copy]
\label{prop:single_prime_copy}
\emph{[PROVED HERE]} For a prime $p$ and $t<p^2$,
\begin{equation}
\label{eq:single-prime-copy}
  \Delta_p(t):=\sum_{a\le t/p}\mu(ap)=-M(t/p).
\end{equation}
Only $p>\sqrt t$ is used; the condition $p>u^2$ of version 5 plays no role.
\end{proposition}

\begin{proof}
$a\le t/p<p$, so $p\nmid a$ and $\mu(ap)=-\mu(a)$.
\end{proof}

\begin{proposition}[Sign content of one prime cone]
\label{prop:prime-cone-signs}
\emph{[PROVED HERE]} Let $\Sq(y)=\sum_{a\le y}\mu(a)^2$. Among $ap$,
$a\le y<p$, the numbers of values $+1,-1,0$ of $\mu$ are
\begin{equation}
\label{eq:cone_counts}
  N_+=\frac{\Sq(y)-M(y)}{2},\qquad N_-=\frac{\Sq(y)+M(y)}{2},\qquad
  N_0=\lfloor y\rfloor-\Sq(y),
\end{equation}
with limiting proportions $3/\pi^2,\ 3/\pi^2,\ 1-6/\pi^2$.
\end{proposition}

\begin{proposition}[Linear window]
\label{prop:linear-new-primes}
\emph{[PROVED HERE]} If $u^2<t<u^4$, no $n\le t$ has two prime factors
$>u^2$ (counted with multiplicity), and
\begin{equation}
\label{eq:first-order-total}
  E_u(t)=-\sum_{u^2<p\le t}M(t/p).
\end{equation}
\end{proposition}

\begin{proposition}[Prime-count form]
\label{prop:prime_count_form}
\emph{[PROVED HERE]} For $u^2<t<u^4$ and $R:=t/u^2$,
\begin{equation}
\label{eq:prime_count_form}
  E_u(t)=-\sum_{a\le R}\mu(a)\bigl(\pi(t/a)-\pi(u^2)\bigr),
  \qquad
  M(t)=S_u(t)-\sum_{a\le R}\mu(a)\pi(t/a)+M(R)\pi(u^2).
\end{equation}
\end{proposition}

\begin{proof}
Expand \eqref{eq:first-order-total} as
$-\sum_{p}\sum_{a\le t/p}\mu(a)$ and interchange: for fixed $a$ the prime
runs over $u^2<p\le t/a$, which is non-empty only if $a<R$.
\end{proof}

\begin{remark}[How the primes above $u^2$ enter]
\label{rem:unknown_primes_enter}
The only input from primes $>u^2$ is the value of $\pi$ at the points $t/a$,
$a\le R$, weighted by $\mu(a)$ at the small scale $R<u^2$. Individual
positions of those primes matter only through these counts.
\end{remark}

\begin{definition}[Copy multiplicities]
\label{def:copy-multiplicity}
For $u^2<t<u^4$, $R=t/u^2$ and $1\le y\le\lfloor R\rfloor$ put
\begin{equation}
\label{eq:Ny-exact}
  N_y(t,u)=\#\Bigl\{p:\ \max\bigl(u^2,\tfrac{t}{y+1}\bigr)<p\le\tfrac ty\Bigr\},
  \qquad
  E_u(t)=-\sum_{1\le y\le\lfloor R\rfloor}N_y(t,u)M(y).
\end{equation}
\end{definition}

\begin{remark}[Asymptotics of the multiplicities, corrected]
\label{rem:copy-count-asymptotic}
For fixed $R>1$ and $u\to\infty$,
\begin{equation}
\label{eq:Ny-fixed-R}
  N_y\sim\frac{Ru^2}{2\log u}\cdot\frac1{y(y+1)}\quad(1\le y\le\lfloor R\rfloor-1),
  \qquad
  N_{\lfloor R\rfloor}\sim\frac{u^2}{2\log u}\Bigl(\frac R{\lfloor R\rfloor}-1\Bigr)
  \quad(R\notin\mathbb{Z}),
\end{equation}
and $N_R=0$ if $R\in\mathbb{Z}$.
\end{remark}

\begin{warning}
\label{warn:v5_Ny}
Version 5 stated the first relation of \eqref{eq:Ny-fixed-R} for all $y<R$.
This is \emph{incorrect} for $y=\lfloor R\rfloor$ when $R\notin\mathbb{Z}$,
since then $t/(y+1)<u^2$.
\end{warning}

\begin{definition}[Coherent coefficient]
\label{def:coherent_coefficient}
\begin{equation}
\label{eq:C_def}
  \Cc(R):=\int_1^R\frac{M(w)}{w^2}\,dw
  =\sum_{a\le R}\mu(a)\Bigl(\frac1a-\frac1R\Bigr)
  =\sum_{a\le R}\frac{\mu(a)}a-\frac{M(R)}R .
\end{equation}
\end{definition}

\begin{proposition}[Fixed $R$, corrected]
\label{prop:fixed_R_corrected}
\emph{[PROVED HERE]} For fixed $R>1$, as $u\to\infty$,
\begin{equation}
\label{eq:E_fixed_R}
  E_u(Ru^2)=-\frac{Ru^2}{2\log u}\bigl(\Cc(R)+o(1)\bigr).
\end{equation}
\end{proposition}

\begin{proof}
In \eqref{eq:prime_count_form} the sum is finite; for $a<R$ the prime number
theorem gives $\pi(Ru^2/a)-\pi(u^2)\sim(R/a-1)u^2/(2\log u)$, and
$\sum_{a\le R}\mu(a)(R/a-1)=R\,\Cc(R)$.
\end{proof}

\begin{lemma}[The coefficient of version 5]
\label{lem:copy-weight-coefficient}
\emph{[PROVED HERE]} The coefficient $A(R)$ of version 5, here
\begin{equation}
\label{eq:Theta_def}
  \Theta(R):=\sum_{1\le y\le R}\frac{M(y)}{y(y+1)}
  =\sum_{n\le R}\frac{\mu(n)}n-\frac{M(R)}{\lfloor R\rfloor+1},
\end{equation}
satisfies
\begin{equation}
\label{eq:Theta_C}
  \Theta(R)=\Cc\bigl(\lfloor R\rfloor+1\bigr),\qquad
  \Cc(R)=\Theta(R)-M(R)\Bigl(\frac1R-\frac1{\lfloor R\rfloor+1}\Bigr),
\end{equation}
and
\begin{equation}
\label{eq:C_tail}
  \Cc(R)=-\int_R^\infty\frac{M(w)}{w^2}\,dw,\qquad
  \Theta(R)=-\sum_{n>R}\frac{M(n)}{n(n+1)}.
\end{equation}
Version 5 wrote $M(R)/(R+1)$ in \eqref{eq:Theta_def}; this is correct only
for integer $R$.
\end{lemma}

\begin{proof}
$M$ is constant on $[y,y+1)$ and $\int_y^{y+1}w^{-2}dw=1/(y(y+1))$. For $s>1$, \eqref{eq:mellin_M} gives
\[
\int_1^\infty M(w)w^{-s-1}\,dw=\frac1{s\zeta(s)}\longrightarrow0
\qquad(s\downarrow1).
\]
For $s$ in a fixed interval $[1,1+\delta]$, the classical bound
$M(w)\ll we^{-c\sqrt{\log w}}$ dominates the integrand by an integrable
multiple of $e^{-c\sqrt{\log w}}/w$.  Dominated convergence therefore gives
$\int_1^\infty M(w)w^{-2}dw=0$.
\end{proof}

\begin{warning}[Main term of version 5]
\label{warn:v5_main_term}
Version 5 asserted $E_u^{(1)}(Ru^2)\sim-\frac{Ru^2}{2\log u}\Theta(R)$. By
\eqref{eq:E_fixed_R} and \eqref{eq:Theta_C} this is \emph{correct only when
$M(\lfloor R\rfloor)=0$}. For $1<R<2$ one has
$E_u=-(\pi(t)-\pi(u^2))\sim-(R-1)\frac{u^2}{2\log u}$, whereas version 5
gives $-\frac R2\cdot\frac{u^2}{2\log u}$.
\end{warning}

\begin{corollary}[Decay of the coefficients]
\label{cor:AR-VK}
\emph{[PROVED HERE]} By \eqref{eq:C_tail} and \eqref{eq:VK_M},
\begin{equation}
\label{eq:AR-VK}
  \Cc(R),\ \Theta(R)\ll\exp\bigl(-c_1(\log R)^{3/5}(\log\log R)^{-1/5}\bigr).
\end{equation}
\end{corollary}

\begin{proposition}[Both parts are large]
\label{prop:both_parts_large}
\emph{[PROVED HERE]}
\begin{enumerate}[label=\textup{(\roman*)}]
\item For fixed $R>1$ with $\Cc(R)\ne0$,
\begin{equation}
\label{eq:E_fixed_R_size}
  E_u(Ru^2)\sim-\frac{Ru^2}{2\log u}\Cc(R),\qquad
  S_u(Ru^2)\sim+\frac{Ru^2}{2\log u}\Cc(R);
\end{equation}
for $1<R<2$, $\Cc(R)=1-1/R>0$.
\item There is $C>0$ such that, as $t\to\infty$, uniformly for $u^2<t<u^4$
with $R=t/u^2\ge\exp(C(\log\log t)^2)$,
\begin{equation}
\label{eq:E_log2}
  E_u(t)=(1+o(1))\frac{t}{(\log t)^2},\qquad
  S_u(t)=-(1+o(1))\frac{t}{(\log t)^2}.
\end{equation}
\end{enumerate}
In regime (i) this comparison is asserted only under the displayed condition
$\Cc(R)\ne0$; in regime (ii), $|E_u(t)|$ and $|S_u(t)|$ exceed
$t^{1/2+\eps}$ for every fixed $\eps<1/2$ and all large $t$.
\end{proposition}

\begin{proof}
We deliberately use the classical weaker exponential form recorded in
\cite{MontgomeryVaughan2007}; the stronger Vinogradov--Korobov form is not
needed here.  For some $c>0$,
\begin{equation}
\label{eq:pnt_inputs}
  \pi(y)=\li(y)+O\bigl(ye^{-c\sqrt{\log y}}\bigr),\quad
  M(y)\ll ye^{-c\sqrt{\log y}},\quad
  \mt(y):=\sum_{a\le y}\frac{\mu(a)}a\ll e^{-c\sqrt{\log y}},
\end{equation}
and, in the usual convergent limiting sense,
$\sum_{n\ge1}\mu(n)\log n/n=-1$ (the value at $s=1$ obtained by
differentiating $1/\zeta(s)$).  This series is \emph{not} asserted to be
absolutely convergent.  In contrast, the integral below is absolutely
convergent because of the displayed bound on $\mt$. Hence
\begin{equation}
\label{eq:mtilde_integral}
  \int_1^\infty\mt(w)\frac{dw}w
  =\lim_{Y\to\infty}\Bigl(\mt(Y)\log Y-\sum_{a\le Y}\frac{\mu(a)\log a}a\Bigr)=1,
\end{equation}
with absolute convergence.

(i) Proposition~\ref{prop:fixed_R_corrected}, and $M(t)=o(t/\log t)$. For
$1<R<2$, $M=1$ on $[1,2)$.

(ii) Since $t<u^4$, $R<\sqrt t$, so $\log(t/a)\ge\frac12\log t$ for
$a\le R$. By \eqref{eq:prime_count_form},
\begin{equation}
\label{eq:E_three_terms}
  E_u(t)=-\sum_{a\le R}\mu(a)\li(t/a)
  -\sum_{a\le R}\mu(a)\bigl(\pi(t/a)-\li(t/a)\bigr)+M(R)\pi(u^2).
\end{equation}
The second term is
$\ll t(1+\log R)e^{-c\sqrt{\log t/2}}=o(t/\log^2t)$, using
$\sum_{a\le R}1/a\ll1+\log R$.  For the third, $t<u^4$ and $t>u^2$ imply
$\tfrac14\log t<\log u<\tfrac12\log t$, so
\[
|M(R)|\pi(u^2)\ll Re^{-c\sqrt{\log R}}\frac{u^2}{\log u}
\ll e^{-c\sqrt{\log R}}\frac{t}{\log t}
=o(t/\log^2t)
\]
once $C$ is chosen so that $c\sqrt C>1$. For the first, with
$F(w)=w\li(t/w)$,
\begin{align}
\label{eq:partial_summation_li}
  \sum_{a\le R}\mu(a)\li(t/a)
  &=\mt(R)F(R)-\int_1^R\mt(w)F'(w)\,dw,\\
  F'(w)
  &=\li(t/w)-\frac{t/w}{\log(t/w)}\notag\\
  &=\frac{t/w}{\log^2(t/w)}
    \Bigl(1+O\Bigl(\frac1{\log t}\Bigr)\Bigr).\notag
\end{align}
The boundary term is $\ll e^{-c\sqrt{\log R}}t/\log t=o(t/\log^2t)$. The integral equals
\[
\frac t{\log^2t}\left(\int_1^R\mt(w)g_t(w)\frac{dw}w+O(1/\log t)\right),
\qquad
g_t(w)=\frac{\log^2t}{\log^2(t/w)}\in[1,4].
\]
To justify the moving endpoint uniformly, extend
$g_t(w)\mathbf1_{\{w\le R\}}$ to $[1,\infty)$.  The hypothesis gives
\[
R\ge\exp\!\bigl(C(\log\log t)^2\bigr)\longrightarrow\infty.
\]
Thus the extended factor tends to $1$ for each fixed $w$, is bounded by $4$,
and $\int_1^\infty|\mt(w)|\,dw/w<\infty$.  Dominated convergence therefore gives
\[
\int_1^R\mt(w)g_t(w)\frac{dw}w=1+o(1)
\]
uniformly in the stated range.  Hence
$E_u(t)=(1+o(1))t/\log^2t$, and $S_u=M-E_u$ with $M(t)=o(t/\log^2t)$.
\end{proof}

\begin{remark}[Why $t/\log^2t$]
\label{rem:why_log2}
$\Cc(R)$ captures only the $1/\log t$ term, which tends to zero. The
$1/\log^2t$ term comes from $\zeta'(s)/\zeta(s)^2\to-1$, that is, from the
pole of $\zeta$ at $s=1$, and it does not decay in $R$. The version 5 claim
that the ``coherent fraction tends to zero'' is true for the first term and
does not describe the size of $E_u$.
\end{remark}

\begin{definition}[Remaining margin, version 5]
\label{def:remaining-margin}
For a target $T(t)=t^{1/2+\eps}$ put
\begin{equation}
\label{eq:margin-def}
  G_{u,\eps}(t)=\max\bigl(0,\ T(t)-|S_u(t)|\bigr),
\end{equation}
and call the large-prime part \emph{causally responsible} for a crossing of
$T$ if $|E_u(t)|\ge G_{u,\eps}(t)$ at the crossing.
\end{definition}

\begin{warning}[The margin is vacuous]
\label{warn:margin_vacuous}
With $S_u$ as in \eqref{eq:baseline_channel_def},
Proposition~\ref{prop:both_parts_large} gives $|S_u(t)|>T(t)$, hence
$G_{u,\eps}(t)=0$, for all large $t$ in both regimes. The test is always
passed and does not discriminate. ``Causal responsibility'' of one of two
cancelling terms of size $t/\log^kt$ for the value of their sum is not well
defined.
\end{warning}

\begin{warning}[Calibration remarks of version 5]
\label{warn:v5_calibration}
Version 5 required $|A(R)|\ll t^{-1/2+\eps}\log t$ for the large-prime part
to stay at the RH scale, and for $t=u^\alpha$ derived
$|A(R)|\ll R^{-\eta(\alpha,\eps)}$ with
$\eta=\alpha(1/2-\eps)/(\alpha-2)$ ($\eta=3/2,1,3/4,5/8$ at
$\alpha=3,4,6,10$). These remarks are \emph{superseded}: the coefficient is
$\Cc$, not $\Theta$ (Warning~\ref{warn:v5_main_term}); for $\alpha>4$ the
point $t$ lies outside the window of
Proposition~\ref{prop:linear-new-primes}, where
\eqref{eq:single-prime-copy} fails for $u^2<p\le\sqrt t$; and by
Proposition~\ref{prop:both_parts_large}, $E_u$ alone is never at the RH
scale, so only $S_u+E_u=M$ can be compared with $T$.
\end{warning}

\begin{remark}[First copy transitions]
\label{rem:first-copy-transitions}
For $1<R<2$ every copy has length $1$ and
$E_u(t)=-(\pi(t)-\pi(u^2))$, consistent with $\Cc(R)=1-1/R$. For $2<R<3$
copies of length $2$ appear but $M(2)=0$ ($\Cc=1/2$). For $3<R<4$ copies of
length $3$ enter with $M(3)=-1$, of opposite net sign.
\end{remark}

\section{Local large-prime decomposition of an excursion}
\label{sec:local_large_prime}

\noindent\textbf{Status: MIXED.} Proposition~\ref{prop:local_large_prime}
and Lemma~\ref{lem:inverse_log_sum} are [PROVED HERE]; the decomposition has
classical sieve/Buchstab ancestry. The input on
primes in short intervals is from the literature \cite{Huxley1972};
Heuristic~\ref{heu:local_drift} is heuristic.

\begin{proposition}[Local decomposition]
\label{prop:local_large_prime}
\emph{[PROVED HERE]} Let $I_\nu=(u,v]$ be a complete excursion and
$y=\sqrt v$. For $u<t\le v$,
\begin{equation}
\label{eq:local_split}
  M(t)=\sum_{\substack{u<n\le t\\ \Pp(n)\le y}}\mu(n)
  -\sum_{a<y}\mu(a)\,\Pi_a(t),\qquad
  \Pi_a(t):=\#\Bigl\{p:\ \max\bigl(y,\tfrac ua\bigr)<p\le\tfrac ta\Bigr\},
\end{equation}
and in particular the bridge condition $M(v)=0$ reads
\begin{equation}
\label{eq:local_bridge}
  \sum_{\substack{u<n\le v\\ \Pp(n)\le\sqrt v}}\mu(n)
  =\sum_{a<\sqrt v}\mu(a)\,\Pi_a(v).
\end{equation}
\end{proposition}

\begin{proof}
$M(u)=0$. An $n\le v=y^2$ has at most one prime factor $p>y$; then
$n=ap$ with $a<v/y=y<p$, so $\mu(n)=-\mu(a)$, and $(a,p)\mapsto ap$ is a
bijection onto these $n$.
\end{proof}

\begin{remark}[What the large primes do inside an excursion]
\label{rem:local_meaning}
The decomposition holds for every excursion, with threshold $\sqrt{v_\nu}$
and no horizon $u_\nu^2$. The large-prime input consists of prime counts in
the intervals $(u/a,t/a]$, of length at most $\ell_\nu/a$ at height about
$v_\nu/a$. Historically Huxley proved the prime number theorem uniformly in all
intervals of length $h\ge x^{7/12+\varepsilon}$ \cite{Huxley1972}. This is
no longer the record: the zero-density estimates of Guth--Maynard imply the
prime number theorem in all intervals of length $x^{17/30+o(1)}$
\cite{GuthMaynard2024}. For M\"obius sums, Matom\"aki--Ter\"av\"ainen prove
$\sum_{x<n\le x+x^\theta}\mu(n)=o(x^\theta)$ for every
$\theta>0.55$ \cite{MatomakiTeravainen2023}. As of September 2026 these are the exponents recorded here; below those
ranges this paragraph asserts \emph{[NO USEFUL BOUND]} of PNT-asymptotic
strength.
\end{remark}

\begin{lemma}[Inverse logarithms]
\label{lem:inverse_log_sum}
\emph{[PROVED HERE]} With $C$ as in
Proposition~\ref{prop:both_parts_large}, the following holds uniformly for
\[
\exp(C(\log\log x)^2)\le z\le\sqrt x:
\]
\begin{equation}
\label{eq:inverse_log_sum}
  \sum_{a\le z}\frac{\mu(a)}{a\log(x/a)}=-\frac{1+o(1)}{(\log x)^2}.
\end{equation}
\end{lemma}

\begin{proof}
Partial summation gives
\[
\frac{\mt(z)}{\log(x/z)}-
\int_1^z\frac{\mt(w)}{w\log^2(x/w)}\,dw.
\]
Because $z\le\sqrt x$, the boundary term is
$O(e^{-c\sqrt{\log z}}/\log x)=o(1/\log^2x)$ uniformly under the stated
lower bound for $z$.  Write
\[
\frac1{\log^2(x/w)}=\frac1{\log^2x}\,g_x(w),\qquad
1\le g_x(w)\le4.
\]
Extend $g_x(w)\mathbf1_{\{w\le z\}}$ by zero to $[1,\infty)$.  Since
$z\to\infty$ uniformly, this factor tends pointwise to $1$, while
$|\mt(w)|/w$ is integrable by \eqref{eq:mtilde_integral}.  Dominated
convergence gives
\[
\int_1^z\frac{\mt(w)}{w\log^2(x/w)}\,dw
=\frac{1+o(1)}{\log^2x},
\]
which proves \eqref{eq:inverse_log_sum} with the displayed minus sign.
\end{proof}

\begin{heuristic}[Opposite drifts]
\label{heu:local_drift}
This is a heuristic argument. Let $u\ge\sqrt v\exp(C(\log\log v)^2)$. Near
$t\in I_\nu$ the counts $\Pi_a$ grow for $a<t/\sqrt v$, at rate
$1/(a\log(t/a))$ if primes near $t/a$ have their expected local density
uniformly in $a$ (beyond what is known). Ignoring these difficulties, over
$(t,t+h]$ the large-prime part changes by about
$-h\sum_{a<t/\sqrt v}\mu(a)/(a\log(t/a))=(1+o(1))h/\log^2t$ by
Lemma~\ref{lem:inverse_log_sum}. Since $M(v)-M(u)=0$, the
$\sqrt v$-smooth part drifts at rate about $-1/\log^2v$. The drift does not
depend on $\sigma_\nu$. Indeed,
\[
\frac{\ell_\nu}{\log^2v_\nu}\gg\sqrt{\ell_\nu}
\quad\Longleftrightarrow\quad
\ell_\nu\gg\log^4v_\nu.
\]
Thus for $\ell_\nu\gg(\log v_\nu)^4$ both heuristic drift contributions
exceed $\sqrt{\ell_\nu}$, and \eqref{eq:local_bridge} is a balance between two
opposite drifts: the local analogue of
Proposition~\ref{prop:both_parts_large}.
\end{heuristic}

\section{A model for the black curve}
\label{sec:model}

\noindent\emph{Status.} Heuristic throughout.

\begin{heuristic}[Scale and period]
\label{heu:model}
This is a heuristic argument. Motivated by \eqref{eq:explicitM}, one is led
to model $M(x)\approx\sqrt x\,F(\log x)$ with $F$ almost periodic.  As a one-frequency toy model only, retain the
lowest zero $\gamma_1=14.1347\ldots$; RH alone does not imply that its
coefficient dominates the higher frequencies. Within an excursion of length
$\asymp x$ this suggests $H_\nu\asymp\sqrt x$ and
$\sqrt{H_\nu}\ell_\nu\asymp x^{5/4}$, so the black curve should oscillate at
scale
\begin{equation}
\label{eq:model_scale}
  |\Anet(x)|\asymp x^{5/4},
\end{equation}
with a period $2\pi/\gamma_1$ in $\log x$, that is, a factor
$e^{2\pi/\gamma_1}\approx1.56$ in $x$.
\end{heuristic}

\section{Numerical data}
\label{sec:data}

\noindent\emph{Status.} Observations inherited from version 5.  No new
numerical computation is claimed for versions 6--9; version 9 only corrects
the interpretation of the recorded heuristic prediction.

\begin{experiment}[Excursion census]
\label{exp:census}
\emph{Object:} the excursion data $(u_\nu,\tau_\nu,v_\nu,H_\nu,\sigma_\nu)$
and the curves $\Ap,\Am,\Anet$.
\emph{Range:} $n\le N=455\,052\,511$, with $Z(N)=87\,385$ excursions.
\emph{Expected outcome:} the bound \eqref{eq:transfer_bound}.  If one takes
Heuristic~\ref{heu:renewal} literally, then for $L=\sqrt N$ it predicts
\[
\#\{\ell_\nu>\sqrt N\}\asymp N^{1/4},
\]
not $Z(N)N^{-1/4}$.  For the recorded value
$N=455\,052\,511$, $N^{1/4}$ is of order $1.5\times10^2$; the former
``approximately 600'' prediction is withdrawn.
\emph{Observed outcome (version 5):}
\begin{equation}
\label{eq:census_values}
  \Ap(N)=6.3449\times10^9,\quad \Am(N)=-7.0877\times10^9,\quad
  \Anet(N)=-7.4279\times10^8,
\end{equation}
\begin{equation}
\label{eq:census_max}
  \max_\nu\frac{H_\nu}{\sqrt{\ell_\nu}}=2.0743,\qquad
  \max_\nu\frac{\beta_\nu}{\ell_\nu^{3/2}}=1.0214 .
\end{equation}
The count of excursions longer than $\sqrt N$ and the maxima restricted to
$\ell_\nu>\sqrt N$ were not recorded in the manuscript.  Reproducible use
of these observations requires the code, the excursion file or checksum,
the exact area-integration convention, and an explicit convention for the
terminal block. At this range $N^{0.05}\approx2.7$, so
factors $x^\eps$ are not separable from constants.
\end{experiment}

\begin{warning}[Constant in the transfer bound]
\label{warn:census_c}
Version 5 compared $|\Anet(N)|$ with \eqref{eq:transfer_bound} using
$c=\max_{10^3\le n\le N}|M(n)|n^{-1/2}$, reporting a bound about $30$ times
larger than $|\Anet(N)|$. That $c$ does not satisfy the hypothesis of
Lemma~\ref{lem:transfer}, which requires all $n\le x^*$; since $M(1)=1$, any
admissible $c$ is $\ge1$. The comparison must either use the full maximum or
split off the range below $10^3$.
\end{warning}

\section{Open questions and future work}
\label{sec:open_questions}

\noindent\textbf{Status: OPEN / HEURISTIC ONLY.} No item in this section
is asserted as a theorem unless it explicitly cites an earlier theorem.

\begin{enumerate}[label=\textup{(\arabic*)}]
\item \textbf{OPEN, RH-EQUIVALENT.} Question~\ref{q:excursion-constraint}:
the uniform bilinear bound on square-ordinary excursion bridges. There is
currently \emph{no proved RH-scale bound} from the bilinear representation
alone. By Corollary~\ref{cor:power-ordinary-RH}, the uniform mathematical bound
itself is equivalent to RH.
\item \textbf{RESOLVED BY PUBLISHED INPUT.}
Corollary~\ref{cor:finiteness_expanding}: square-exceptional excursions are
finite; more generally, for every fixed $A>1$, only finitely many excursions
satisfy $x_1>x_0^A$. The deduction from $M+2$ uses the combined oscillation
size, not a bare sign change.
\item \textbf{OPEN.} Whether \eqref{eq:rh_area} is equivalent to RH.
The former $5/6$ zero-free deduction is \emph{FALSE / WITHDRAWN}.
\item Whether the pruned local condition \eqref{eq:pruned_Q} holds, and a
third bridge constraint (Remark~\ref{rem:conditioning}).
\item Experiment: the number of excursions with $\ell_\nu>\sqrt N$ against the
renewal-scale prediction $\asymp N^{1/4}$, and the maxima of $H_\nu/\sqrt{\ell_\nu}$ and
$\beta_\nu/\ell_\nu^{3/2}$ over those excursions.
\item Experiment: split each excursion by \eqref{eq:local_split} and test
Heuristic~\ref{heu:local_drift}, in particular the independence of the drift
from $\sigma_\nu$.
\item Experiment: frequency of touching runs (Definition~\ref{def:crossing_touching})
and of ballistic excursions (Heuristic~\ref{heu:local_fails}).
\end{enumerate}

\section*{Notebook History}
\label{sec:history}

\noindent\textbf{Version 5.} Single section: definitions, exact formulas,
internal correlations, bounds, criteria, prime horizon, copy propagation,
model, census data.

\medskip
\noindent\textbf{Version 6.} Restructured into labelled sections following
the authoring rules. Corrections: zero runs and touching zeros
(Definitions~\ref{def:excursion}, \ref{def:crossing_touching},
Lemma~\ref{lem:critical}); Remark~\ref{rem:omega};
Remark~\ref{rem:onedirection}; hypothesis of Lemma~\ref{lem:transfer};
constant in Remark~\ref{rem:chowla-scale}; $j=1$ in
Proposition~\ref{prop:double-exp}; \eqref{eq:Ny-fixed-R}; main term
(Proposition~\ref{prop:fixed_R_corrected}, Warning~\ref{warn:v5_main_term}).
New: Warning~\ref{warn:information_language}, Lemma~\ref{lem:rough_part},
Proposition~\ref{prop:prime_count_form},
Proposition~\ref{prop:both_parts_large},
Section~\ref{sec:local_large_prime}, Heuristic~\ref{heu:local_fails},
Corollary~\ref{cor:pruned_local}, proofs of
Propositions~\ref{prop:L1} and~\ref{prop:dipoles}. Archived as superseded:
the calibration remarks (Warning~\ref{warn:v5_calibration}) and the margin
test (Warning~\ref{warn:margin_vacuous}).

\medskip
\noindent\textbf{Version 7.} Full merge of version~6 with the reviewed
changes to Sections~2--4. New exact material: terminal skeleton,
$g$-skeletons, real/Hermitian triangular identity, extreme-lag identity,
exact terminal quadratic decomposition, sharp $(L,H)$ bounds and the $1/12$
fixed-length constant. New central results:
Lemma~\ref{lem:power-lacunary} ($O(\log\log X)$ power-exception count),
Lemma~\ref{lem:linear-lacunary} ($O(\log X)$ linear-long count), the Kaczorowski--Pintz large-oscillation input, and the unconditional localized
RH criteria, in particular \eqref{eq:square-horizon-RH}. The square range is
written as the exact bilinear form \eqref{eq:bilinear-excursion};
\eqref{eq:HB-k} records why prefix determinacy itself is not special to the
square.

\medskip
\noindent\textbf{Version 8.} Merged the version-7 change block with all of
version 6, added explicit evidentiary statuses, updated Hurst and the short-
interval literature, withdrew the false $5/6$ zero-free inference, and
isolated the sparse-exception localization principle.  A subsequent audit
identified remaining notation, proof-detail and numerical issues; version 8
should therefore be regarded as a transitional merged draft rather than an
audited edition.

\medskip
\noindent\textbf{Version 9.} Revised merged edition.  Main corrections:
separate [PUBLISHED INPUT]/[PROVED HERE]/[CONDITIONAL]/[HEURISTIC] statuses;
unified terminal notation; expanded the triangular and fixed-$(L,H)$ proofs;
made the explicit $\widetilde A$ formula explicitly formal and justified the
residue $12$; corrected the pruning quantifiers and terminal indicators;
added the square-truncated RH criterion; transcribed the Kaczorowski--Pintz
input with page and equation numbers and used their local relations
(2.2)--(2.3) to prove finiteness of the $A$-exceptional Mertens excursions
for every fixed $A>1$; separated the bilinear form
$\mathcal B_u$ from the smooth baseline $S_u$; supplied the order-$k$
truncation proof; completed the moving-endpoint dominated-convergence steps
in the large-prime and inverse-logarithm arguments; replaced the unsupported
``lowest-zero dominance'' by a one-frequency toy model; and withdrew the
incorrect numerical prediction $Z(N)N^{-1/4}\approx600$.

\begin{thebibliography}{99}


\bibitem{CP17}
P.~Carmona and N.~P\'etr\'elis, \emph{Limit theorems for random walk excursion
conditioned to enclose a typical area}, arXiv:1709.06448.

\bibitem{DKW15}
D.~Denisov, M.~Kolb and V.~Wachtel, \emph{Local asymptotics for the area of
random walk excursions}, J. London Math. Soc. \textbf{91} (2015), 495--513.

\bibitem{HB82}
D.~R. Heath-Brown, \emph{Prime numbers in short intervals and a generalized
Vaughan identity}, Canad. J. Math. \textbf{34} (1982), 1365--1377.

\bibitem{HT23}
H.~A. Helfgott and L.~Thompson, \emph{Summing $\mu(n)$: a faster elementary
algorithm}, Res. Number Theory \textbf{9} (2023), Paper No.~6.

\bibitem{GuthMaynard2024}
L.~Guth and J.~Maynard, \emph{New large value estimates for Dirichlet
polynomials}, Ann. of Math. (2) \textbf{203} (2026), no.~2, 623--675,
doi:10.4007/annals.2026.203.2.6. Corollary~1.3 gives the prime number
theorem in intervals of length $x^{17/30+o(1)}$.

\bibitem{MatomakiTeravainen2023}
K.~Matom\"aki and J.~Ter\"av\"ainen, \emph{On the M\"obius function in
all short intervals}, J. Eur. Math. Soc. \textbf{25} (2023), no.~4,
1207--1225. doi:10.4171/JEMS/1205.

\bibitem{Hurst18}
G.~Hurst, \emph{Computations of the Mertens function and improved bounds on
the Mertens conjecture}, Math. Comp. \textbf{87} (2018), no.~310, 1013--1028.

\bibitem{Janson07}
S.~Janson, \emph{Brownian excursion area, Wright's constants in graph
enumeration, and other Brownian areas}, Probab. Surveys \textbf{4} (2007),
80--145.

\bibitem{KP86}
J.~Kaczorowski and J.~Pintz, \emph{Oscillatory properties of arithmetical
functions. I}, Acta Math. Hungar. \textbf{48} (1986), 173--185.

\bibitem{NguyenThe04}
M.~Nguy\^en Th\^e, \emph{Area and inertial moment of Dyck paths}, Combin.
Probab. Comput. \textbf{13} (2004), 697--716.

\bibitem{Richard05}
C.~Richard, \emph{On $q$-functional equations and excursion moments},
arXiv:math/0503198.

\bibitem{Takacs91}
L.~Tak\'acs, \emph{A Bernoulli excursion and its various applications},
Adv. Appl. Probab. \textbf{23} (1991), 557--585.

\bibitem{Huxley1972}
M.~N. Huxley, On the difference between consecutive primes,
\emph{Invent. Math.} \textbf{15} (1972), 164--170.
doi:10.1007/BF01418933.

\bibitem{KotnikTeRiele2006}
T.~Kotnik and H.~te Riele, The Mertens conjecture revisited, in
\emph{Algorithmic Number Theory (ANTS VII)}, Lecture Notes in Comput. Sci.
\textbf{4076}, Springer, Berlin, 2006. doi:10.1007/11792086\_12.

\bibitem{MontgomeryVaughan2007}
H.~L. Montgomery and R.~C. Vaughan, \emph{Multiplicative Number Theory I:
Classical Theory}, Cambridge Stud. Adv. Math. \textbf{97}, Cambridge
University Press, 2007. ISBN 978-0-521-84903-6.

\bibitem{Ng2004}
N.~Ng, The distribution of the summatory function of the M\"obius function,
\emph{Proc. London Math. Soc. (3)} \textbf{89} (2004), no.~2, 361--389.
doi:10.1112/S0024611504014741. arXiv:math/0310381.

\bibitem{OdlyzkoTeRiele1985}
A.~M. Odlyzko and H.~J.~J. te Riele, Disproof of the Mertens conjecture,
\emph{J. Reine Angew. Math.} \textbf{357} (1985), 138--160. MR0783538.

\bibitem{Titchmarsh1986}
E.~C. Titchmarsh, \emph{The Theory of the Riemann Zeta-Function}, 2nd ed.,
revised by D.~R. Heath-Brown, Clarendon Press, Oxford, 1986.

\bibitem{Walfisz1963}
A.~Walfisz, \emph{Weylsche Exponentialsummen in der neueren Zahlentheorie},
VEB Deutscher Verlag der Wissenschaften, Berlin, 1963.

\end{thebibliography}

\end{document}
